\documentclass[11pt]{article}

\usepackage[margin=1in]{geometry}
\usepackage[T1]{fontenc}
\usepackage[utf8]{inputenc}
\usepackage{lmodern}
\usepackage{microtype}
\usepackage{amsmath,amssymb,amsthm,mathtools}
\usepackage{booktabs}
\usepackage{enumitem}
\usepackage{xcolor}
\usepackage{hyperref}
\usepackage{cleveref}
\usepackage{natbib}

\hypersetup{
  colorlinks=true,
  linkcolor=blue!70!black,
  citecolor=green!40!black,
  urlcolor=blue!70!black
}

\newtheorem{proposition}{Proposition}
\newtheorem{definition}{Definition}
\newtheorem{remark}{Remark}

\newcommand{\E}{\mathbb{E}}
\newcommand{\R}{\mathbb{R}}
\newcommand{\1}{\mathbf{1}}
\DeclareMathOperator*{\argmax}{arg\,max}
\DeclareMathOperator*{\argmin}{arg\,min}

\title{Solving a Risky-Debt Corporate Finance Model with Deep Learning\\
{\large a toy example on the possibility of applying machine learning to finance problems}}
\author{Chak Wong}
\date{\today}

\begin{document}
\maketitle

\begin{abstract}
These notes explain how to adapt the deep-learning method of
\citet{maliar2021deep} to the risky-debt dynamic corporate finance model in
Section 3.6 of \citet{strebulaev2012dynamic}. The goal is not to prove that a
particular neural network will always converge. The goal is more modest and
more useful for computation: we rewrite the firm's recursive problem, the
limited-liability default condition, and the lender's zero-profit debt-pricing
condition as residual equations that can be minimized by stochastic gradient
descent. The exposition is written for early undergraduates. It uses functions,
sums, informal derivatives, expected values, and a small amount of optimization
language; each specialized idea is introduced before it is used.
\end{abstract}

\tableofcontents

\section{What problem are we solving?}

Dynamic corporate finance models describe a firm that makes decisions today,
then faces uncertainty tomorrow. In the simple investment model of
\citet[Section 3.1]{strebulaev2012dynamic}, the firm chooses next period's
capital stock. In the risky-debt model of
\citet[Section 3.6]{strebulaev2012dynamic}, the firm chooses both next period's
capital and next period's debt, and debt is risky because the firm may default.

The difficult object is the \emph{value function}. A value function is a rule
that assigns a number to each state of the firm. For example,
\[
  V(k,b,z)
\]
means the equity value of a firm with capital \(k\), debt \(b\), and productivity
shock \(z\). The state variables are:
\begin{itemize}[leftmargin=2em]
  \item \(k\): beginning-of-period capital;
  \item \(b\): debt due at the beginning of the period, with negative \(b\)
  interpreted as cash;
  \item \(z\): a productivity or profitability shock.
\end{itemize}

\citet[Definition 2.8 and Algorithm 1]{maliar2021deep} show how a dynamic
economic model can be turned into a residual-minimization problem. The broad
idea is:
\begin{enumerate}[leftmargin=2em]
  \item Write the economic model as equations.
  \item Approximate unknown functions, such as \(V\), by neural networks.
  \item Measure how badly the neural networks violate the model equations.
  \item Train the networks to make those violations small.
\end{enumerate}

The central question of this note is: can we do this for the risky-debt model?
The answer is yes, with one important qualification. We must not minimize only
the Bellman equation. We must also impose the lender's debt-pricing equation,
because the risky interest rate is itself an equilibrium object.

\section{Expected values and the shock process}

The firm does not know next period's shock \(z'\) when it makes today's choices.
Following \citet[eq. (3.9)]{strebulaev2012dynamic}, assume that the shock follows
the log autoregression
\begin{equation}
  \log z' = \rho \log z + \varepsilon',
  \label{eq:ar1-shock}
\end{equation}
where \(\rho\) is a persistence parameter and \(\varepsilon'\) is a random
innovation. If we know today's \(z\), then \cref{eq:ar1-shock} gives the
distribution of tomorrow's \(z'\).

If \(F(z')\) is any function of tomorrow's shock, then
\[
  \E[F(z')\mid z]
\]
means the average value of \(F(z')\), using the probabilities implied by today's
state \(z\). If the shock has finitely many possible next values
\[
  z'_1,\ldots,z'_n
\]
with probabilities \(p_1,\ldots,p_n\), then
\[
  \E[F(z')\mid z] = p_1F(z'_1)+\cdots+p_nF(z'_n).
\]
In the continuous case, the sum becomes an integral. These notes use expectation
notation because it covers both cases.

\section{Warm-up: the basic investment Bellman equation}

In the basic model of \citet[Section 3.1]{strebulaev2012dynamic}, the firm has
capital \(k\), chooses next period's capital \(k'\), and therefore invests
\begin{equation}
  I = k'-(1-\delta)k,
  \label{eq:investment-law}
\end{equation}
where \(\delta\) is the depreciation rate. The one-period cash flow is
\begin{equation}
  e(k,k',z)
  =
  \pi(k,z)-\psi(k'-(1-\delta)k,k)-\bigl(k'-(1-\delta)k\bigr).
  \label{eq:basic-cash-flow}
\end{equation}
Here \(\pi(k,z)\) is profit and \(\psi(I,k)\) is an adjustment cost.

Let
\[
  \beta = \frac{1}{1+r},
\]
where \(r\) is the risk-free interest rate. The number \(\beta\) is the one-period
discount factor.

\begin{proposition}[Bellman equation for the basic investment model]
\label{prop:basic-bellman}
Suppose the firm maximizes the expected discounted sum of future cash flows.
Then the value function satisfies
\begin{equation}
  V(k,z)
  =
  \max_{k'}
  \left\{
    e(k,k',z)
    +\beta \E[V(k',z')\mid z]
  \right\}.
  \label{eq:basic-bellman}
\end{equation}
\end{proposition}

\begin{proof}
The value of the firm is the value of today's cash flow plus the discounted
value of all future cash flows. If the firm chooses \(k'\), today's cash flow is
\(e(k,k',z)\). Tomorrow the firm will be in state \((k',z')\), so tomorrow's
value is \(V(k',z')\). Because \(z'\) is random, the firm uses the average value
\(\E[V(k',z')\mid z]\). Because tomorrow's value is one period in the future, it
is multiplied by \(\beta\). Finally, the firm chooses the \(k'\) that gives the
largest total. This gives \cref{eq:basic-bellman}.
\end{proof}

This has the same recursive structure as
\citet[eq. (3.10)]{strebulaev2012dynamic}, but written with \(\beta\) instead
of \(1/(1+r)\).

\section{The risky-debt model}

Now add one-period debt. The firm starts the period with debt \(b\), chooses next
period's capital \(k'\), and chooses next period's debt \(b'\). The important new
feature in \citet[Section 3.6]{strebulaev2012dynamic} is that debt is risky. A
lender will charge a risky interest rate
\[
  \widetilde r(z,k',b')
\]
that depends on today's shock \(z\) and on the firm's promised next-period
capital and debt \((k',b')\).

\subsection{Cash flow with risky debt}

\citet[eq. (3.26)]{strebulaev2012dynamic} define the one-period cash flow as
\begin{align}
  e(k,k',b,b',z;\widetilde r)
  &=
  (1-\tau)\pi(k,z)
  -\psi(k'-(1-\delta)k,k)
  -\bigl(k'-(1-\delta)k\bigr) \notag\\
  &\quad
  +\frac{b'}{1+\widetilde r(z,k',b')}
  +\frac{\tau \widetilde r(z,k',b')b'}
  {(1+\widetilde r(z,k',b'))(1+r)}
  -b .
  \label{eq:risky-cash-flow}
\end{align}

The terms have simple interpretations:
\begin{itemize}[leftmargin=2em]
  \item \((1-\tau)\pi(k,z)\) is after-tax profit.
  \item \(\psi(k'-(1-\delta)k,k)\) is the adjustment cost of investment.
  \item \(k'-(1-\delta)k\) is physical investment.
  \item \(b'/(1+\widetilde r)\) is the amount of cash raised today by issuing
  a promise to repay \(b'\) next period.
  \item The term with \(\tau \widetilde r b'\) is the present value of the
  interest tax deduction used in the Section 3.6 model.
  \item The final \(-b\) is repayment of debt already due today.
\end{itemize}

The firm may also face costly external equity finance. To avoid confusing the
cash-flow function \(e(\cdot)\) with the argument of the financing-cost function,
write a generic pre-financing cash flow as \(x\). Let \(\eta(x)\) be the
additional value effect of external financing costs. In many corporate finance
models, \(\eta(x)=0\) when \(x\ge 0\), because the firm is distributing cash, and
\(\eta(x)\le 0\) when \(x<0\), because the firm must raise costly outside equity.
The exact functional form is not needed for the derivation below.

\subsection{Default and recovery}

If the firm defaults next period, the lender recovers
\begin{equation}
  R(k',z')
  =
  (1-\alpha)\bigl((1-\tau)\pi(k',z')+(1-\delta)k'\bigr),
  \label{eq:recovery}
\end{equation}
where \(\alpha\in[0,1]\) measures deadweight default costs. This is the recovery
formula stated in \citet[Section 3.6]{strebulaev2012dynamic}.

Limited liability means equityholders can walk away instead of accepting a
negative equity value. Thus equity value is never below zero.

\begin{remark}[Do not mix up equity value and debt value]
\label{rem:equity-debt-total-value}
This is the point at which the notation can become confusing. In these notes,
\(V(k,b,z)\) is the \emph{equity} value. It is the value to shareholders, not the
value of the lender's debt claim and not the total value of all claims on the
firm. Therefore, if the firm defaults, shareholders receive zero by limited
liability. The recovery \(R(k',z')\) in \cref{eq:recovery} is paid to lenders, so it
belongs in the debt-pricing equation, not as a positive payoff inside the equity
continuation value.

If instead we were writing a value function for lenders, or a total firm value
function that adds equity and debt together, then default recovery would appear
directly in that object's payoff. The Section 3.6 Bellman equation is an equity
Bellman equation, so default gives equityholders zero while recovery affects
equity indirectly through the risky interest rate.
\end{remark}

\begin{definition}[Policy function]
\label{def:policy-function}
A \emph{policy function} is a rule that tells the firm what to do at every
possible state. In the risky-debt model, a stationary Markov policy is a pair of
functions
\[
  h(k,b,z)=\bigl(h^k(k,b,z),h^b(k,b,z)\bigr),
\]
where \(h^k(k,b,z)\) is the next-period capital chosen at state \((k,b,z)\), and
\(h^b(k,b,z)\) is the next-period debt chosen at that state. The word
\emph{stationary} means that the same rule is used in every period. The word
\emph{Markov} means that the rule depends only on the current state
\((k,b,z)\), not on the full past history.
\end{definition}

\begin{definition}[Policy-specific and optimal going-concern values]
\label{def:going-concern-value}
There are two different objects that should not be confused.

First, suppose the firm is committed to a particular policy \(h\). Given a
policy-specific equity value \(V^h\) and a risky rate \(\widetilde r\), define
the value of taking action \((k',b')\) today and then following policy \(h\) from
tomorrow onward by
\begin{equation}
  Q^h(k,b,z;k',b')
  =
  e(k,k',b,b',z;\widetilde r)
  +\eta(e(k,k',b,b',z;\widetilde r))
  +\beta \E[V^h(k',b',z')\mid z].
  \label{eq:policy-q-value}
\end{equation}
If the firm follows the policy \(h\) today, it chooses
\((h^k(k,b,z),h^b(k,b,z))\), so its policy-specific going-concern value is
\begin{equation}
  G^h(k,b,z)
  =
  Q^h(k,b,z;h^k(k,b,z),h^b(k,b,z)).
  \label{eq:policy-going-concern}
\end{equation}
Policy evaluation therefore uses no maximization over actions:
\begin{equation}
  V^h(k,b,z)=\max\{0,G^h(k,b,z)\}.
  \label{eq:policy-bellman}
\end{equation}

Second, if we are solving for the best policy, we use Bellman optimality. Let
\(V^\star\) denote the optimal equity value. Define
\begin{equation}
  Q^\star(k,b,z;k',b')
  =
  e(k,k',b,b',z;\widetilde r)
  +\eta(e(k,k',b,b',z;\widetilde r))
  +\beta \E[V^\star(k',b',z')\mid z].
  \label{eq:q-value}
\end{equation}
This is the optimal action-value expression. It asks: if the firm takes the
particular action \((k',b')\) today, and then behaves optimally from tomorrow
onward, what is today's equity value? The continuation term
\(\E[V^\star(k',b',z')\mid z]\) is already an equity value with limited
liability built into it. It averages next period's shareholder value, which is
zero in future default states. Define the best going-concern equity value by
\begin{equation}
  G^\star(k,b,z) = \max_{k',b'} Q^\star(k,b,z;k',b').
  \label{eq:best-going-concern}
\end{equation}
When the maximum is attained, an optimal policy \(h^\star\) chooses an
action that attains it at every state:
\[
  h^\star(k,b,z)
  \in
  \argmax_{k',b'} Q^\star(k,b,z;k',b').
\]
\end{definition}

\begin{remark}[Connection with dynamic programming and reinforcement learning]
\label{rem:dp-rl-policy}
In reinforcement-learning notation, \(Q^h\) is the action-value function for a
fixed policy, while \(Q^\star\) is the optimal action-value function. Evaluating
an existing policy uses \(Q^h\) and \(V^h\). Improving or solving for the optimal
policy uses the maximization in \(G^\star\). To match the notation in
\citet[Section 3.6]{strebulaev2012dynamic}, the rest of these notes often drops
the star and writes \(V,Q,G,h\) for the optimal objects
\(V^\star,Q^\star,G^\star,h^\star\). Whenever we evaluate a fixed nonoptimal
policy, we keep the superscript \(h\).
\end{remark}

\begin{proposition}[Risky-debt Bellman equation]
\label{prop:risky-bellman}
In the risky-debt model, optimal equity value satisfies
\begin{equation}
  V^\star(k,b,z)
  =
  \max\{0,G^\star(k,b,z)\}.
  \label{eq:risky-bellman}
\end{equation}
\end{proposition}

\begin{proof}
If the firm continues, the best value it can obtain is \(G^\star(k,b,z)\),
because \(G^\star\) is defined as the best equity value over all choices of
\((k',b')\), followed by optimal behavior in the future. If that value is
negative, equityholders prefer to default and receive zero rather than keep a
negative-value equity claim. Therefore equityholders choose the larger of zero
and the optimal going-concern equity value. This gives
\cref{eq:risky-bellman}, matching \citet[eq. (3.28)]{strebulaev2012dynamic}.
The lender's recovery is not missing; it enters the lender's pricing equation
below and changes current equity value through the risky rate
\(\widetilde r(z,k',b')\). It also matters recursively because future equity
values are computed using future debt-pricing equations.
\end{proof}

\section{Debt pricing}

The risky interest rate is not chosen freely by the firm. Competitive lenders
must break even. If the lender gives the firm cash today in exchange for a
promise \(b'\), then the expected payoff next period must equal what the lender
could earn at the risk-free rate.
In this note, as in the Section 3.6 formula, \(b'\) denotes the debt position
chosen today for next period. We use the paper's payoff convention directly:
when the firm is solvent next period, lenders receive \(b'(1+\widetilde r)\);
when the firm defaults, lenders receive the recovery value.

Let
\[
  D(k',b',z')=
  \begin{cases}
    1, & \text{if the firm defaults next period},\\
    0, & \text{if the firm remains solvent next period}.
  \end{cases}
\]
Equivalently, using \cref{eq:risky-bellman}, the next-period state is a default
state when its best going-concern equity value is nonpositive:
\begin{equation}
  D(k',b',z')=\1\{G^\star(k',b',z')\le 0\}.
  \label{eq:default-set}
\end{equation}

\begin{proposition}[Zero-profit risky debt pricing]
\label{prop:risky-pricing}
For \(b'>0\), the zero-profit condition for the risky debt rate is
\begin{equation}
  b'(1+r)
  =
  \E\left[
    D(k',b',z')R(k',z')
    +(1-D(k',b',z'))b'(1+\widetilde r(z,k',b'))
    \mid z
  \right].
  \label{eq:risky-pricing}
\end{equation}
\end{proposition}

\begin{proof}
The left side, \(b'(1+r)\), is the payoff a lender requires next period to earn
the risk-free return on the chosen debt position \(b'\). On the right side,
there are two cases. If the firm defaults, \(D=1\), and the lender receives the
recovery value \(R(k',z')\). If the firm is solvent, then \(D=0\), and the lender
receives the promised payoff \(b'(1+\widetilde r)\). The lender prices debt so
that the expected payoff equals the required risk-free payoff. This is
\citet[eq. (3.27)]{strebulaev2012dynamic} written with an indicator variable.
\end{proof}

\begin{remark}[Solving explicitly for the risky rate]
\label{rem:explicit-rate}
If \(b'>0\) and the probability of solvency is positive, then
\cref{eq:risky-pricing} can be rearranged as
\begin{equation}
  1+\widetilde r(z,k',b')
  =
  \frac{
    b'(1+r)-\E[D(k',b',z')R(k',z')\mid z]
  }{
    b'\E[1-D(k',b',z')\mid z]
  }.
  \label{eq:explicit-rate}
\end{equation}
This formula shows why the risky rate depends on default probabilities. If
default becomes more likely, the denominator falls and the risky rate generally
rises, all else equal. If \(b'\le 0\), the firm is not issuing positive debt in
this simplified pricing equation; a separate cash-return convention should be
specified.
\end{remark}

\begin{proposition}[The risky-debt model is a coupled fixed point]
\label{prop:coupled-fixed-point}
A solution of the risky-debt model is a pair \((V^\star,\widetilde r)\), together
with optimal policies, satisfying all of the following statements at the same
time:
\begin{align}
  V^\star(k,b,z)
  &=
  \max\{0,G^\star(k,b,z)\}, \label{eq:fixed-point-equity}\\
  D(k,b,z)
  &=
  \1\{G^\star(k,b,z)\le 0\}, \label{eq:fixed-point-default}\\
  b'(1+r)
  &=
  \E\left[
    D(k',b',z')R(k',z')
    +(1-D(k',b',z'))b'(1+\widetilde r(z,k',b'))
    \mid z
  \right],
  \label{eq:fixed-point-pricing}
\end{align}
where \(G^\star\) is computed from \cref{eq:best-going-concern} using the same
\(V^\star\) and \(\widetilde r\). The pricing equation is evaluated for each
positive promised debt choice \(b'>0\), in particular for the policy choice made
by the firm.
\end{proposition}

\begin{proof}
\Cref{eq:fixed-point-equity} is the limited-liability equity Bellman equation.
\Cref{eq:fixed-point-default} says that default is chosen exactly when
continuing gives shareholders no positive value. \Cref{eq:fixed-point-pricing}
is the lender's zero-profit condition. These equations must hold together
because each object uses the others: \(V^\star\) determines the default set, the
default set determines the risky rate through lender recovery, and the risky
rate enters the cash flow used to compute \(G^\star\) and therefore
\(V^\star\). This channel operates both today and in future continuation values.
Thus the problem is not a one-way recursion. It is a fixed point in which the
equity value, default rule, optimal policy, and risky interest rate are solved
jointly.
\end{proof}

\section{Writing default as equations}

The expression \(V=\max\{0,G\}\) is intuitive, but neural-network training often
works better when a condition is written as equations. We can rewrite the max
condition using a standard complementarity idea.

\begin{proposition}[Limited liability as complementarity]
\label{prop:limited-liability-complementarity}
For two real numbers \(V\) and \(G\), the statement
\[
  V=\max\{0,G\}
\]
is equivalent to the three conditions
\begin{equation}
  V\ge 0,\qquad V-G\ge 0,\qquad V(V-G)=0.
  \label{eq:complementarity}
\end{equation}
\end{proposition}

\begin{proof}
First suppose \(V=\max\{0,G\}\). If \(G\ge 0\), then \(V=G\), so \(V\ge 0\),
\(V-G=0\), and \(V(V-G)=0\). If \(G<0\), then \(V=0\), so \(V\ge 0\),
\(V-G=-G>0\), and again \(V(V-G)=0\).

Now suppose the three conditions in \cref{eq:complementarity} hold. Since
\(V(V-G)=0\), at least one of \(V\) and \(V-G\) must be zero. If \(V=0\), then
\(V-G\ge 0\) implies \(G\le 0\), so \(V=\max\{0,G\}\). If \(V-G=0\), then
\(V=G\), and \(V\ge 0\), so \(G\ge 0\) and again \(V=\max\{0,G\}\).
\end{proof}

One smooth way to combine the three complementarity conditions is the
Fischer-Burmeister function
\begin{equation}
  \operatorname{FB}(a,c)=\sqrt{a^2+c^2}-a-c.
  \label{eq:fb}
\end{equation}
It equals zero exactly when \(a\ge 0\), \(c\ge 0\), and \(ac=0\). Maliar et al.
use this function for inequality constraints in their Bellman-equation example
\citep[Section 4.5]{maliar2021deep}.

\section{Neural-network approximations}

Now approximate the unknown objects by functions with trainable parameters
\(\theta\). The functions \(\widehat h^k_\theta\) and \(\widehat h^b_\theta\)
approximate the optimal policy functions introduced in
\cref{def:policy-function}. The value network approximates \(V^\star\), not the
value of an arbitrary initial policy.
Write
\begin{align}
  \widehat V_\theta(k,b,z) &\approx V^\star(k,b,z), \label{eq:v-network}\\
  \widehat h^k_\theta(k,b,z) &\approx \text{optimal }k', \label{eq:k-policy-network}\\
  \widehat h^b_\theta(k,b,z) &\approx \text{optimal }b', \label{eq:b-policy-network}\\
  \widehat r_\theta(z,k',b') &\approx \widetilde r(z,k',b').
  \label{eq:rate-network}
\end{align}
The hats remind us that these are approximations.

At a state \(s=(k,b,z)\), define the network-implied controls
\begin{equation}
  \widehat k'=\widehat h^k_\theta(k,b,z),
  \qquad
  \widehat b'=\widehat h^b_\theta(k,b,z).
  \label{eq:network-policies}
\end{equation}

In an implementation, one often uses transformations to enforce constraints.
For example, if capital must satisfy \(\underline k\le k'\le \overline k\), one
can set
\[
  \widehat h^k_\theta(k,b,z)
  =
  \underline k+(\overline k-\underline k)\sigma(N^k_\theta(k,b,z)),
\]
where \(N^k_\theta\) is an unrestricted neural network and
\(\sigma(x)=1/(1+e^{-x})\) maps every real number into \((0,1)\).

\section{Residuals for the risky-debt model}

The model contains three main mathematical requirements:
\begin{enumerate}[leftmargin=2em]
  \item optimal equity value must satisfy limited liability and the Bellman
  optimality equation;
  \item the chosen policy \((\widehat k',\widehat b')\) must maximize the
  going-concern value;
  \item the risky rate must satisfy the lender zero-profit condition.
\end{enumerate}
We now turn each requirement into a residual. A residual is zero when the
equation is satisfied and nonzero when the equation is violated.

\subsection{Bellman and default residual}

Define the network-implied optimal going-concern value at state \(s=(k,b,z)\):
\begin{align}
  \widehat G_\theta(k,b,z)
  &=
  e(k,\widehat k',b,\widehat b',z;\widehat r_\theta)
  +\eta(e(k,\widehat k',b,\widehat b',z;\widehat r_\theta)) \notag\\
  &\quad
  +\beta \E[
    \widehat V_\theta(\widehat k',\widehat b',z')
    \mid z].
  \label{eq:network-going-concern}
\end{align}
Then the limited-liability Bellman residual is
\begin{equation}
  \mathcal B_\theta(k,b,z)
  =
  \operatorname{FB}\left(
    \widehat V_\theta(k,b,z),
    \widehat V_\theta(k,b,z)-\widehat G_\theta(k,b,z)
  \right).
  \label{eq:bellman-fb-residual}
\end{equation}
This residual is zero exactly when
\[
  \widehat V_\theta(k,b,z)=\max\{0,\widehat G_\theta(k,b,z)\}.
\]
For a future state \((k',b',z')\), the expression
\(\widehat G_\theta(k',b',z')\) means the same construction applied at that
future state, using the network's candidate optimal policy choices for the
period after that. It is the network analogue of \(G^\star\), not the value of a
fixed initial policy \(h\).

\subsection{Debt-pricing residual}

Default next period occurs when the next-period optimal equity value is zero:
\[
  V^\star(k',b',z')=0.
\]
Using the optimal Bellman relation \(V^\star=\max\{0,G^\star\}\), this is
equivalent to \(G^\star(k',b',z')\le 0\). Therefore the default indicator implied
by the network is
\begin{equation}
  \widehat D_\theta(k',b',z')
  =
  \1\{\widehat G_\theta(k',b',z')\le 0\}.
  \label{eq:default-indicator}
\end{equation}
This is an implementation choice based on the Section 3.6 default logic: the
paper states that the firm defaults when equity value is zero, while the network
uses the approximated going-concern value to classify that event. In the exact
model, \(V^\star=0\) and \(G^\star\le 0\) are equivalent. For the approximating
functions, this equivalence is exact only when the Bellman residual is zero, and
otherwise it is an approximation whose quality should be checked. For numerical
training, the sharp indicator can be replaced by a smooth approximation, for
example
\[
  \widehat D_{\theta,\kappa}(k',b',z')
  =
  \frac{1}{1+\exp(\kappa \widehat G_\theta(k',b',z'))},
\]
where large \(\kappa>0\) makes the smooth function close to an indicator. With
the sharp indicator, the pricing equation below represents the original default
classification; with the smooth indicator, it represents a smoothed approximate
model.

For positive promised debt \(b'>0\), define the pricing residual
\begin{align}
  \mathcal P_\theta(z,k',b')
  &=
  b'(1+r) \notag\\
  &\quad
  -
  \E\left[
    \widehat D_\theta(k',b',z')R(k',z')
    +(1-\widehat D_\theta(k',b',z'))b'(1+\widehat r_\theta(z,k',b'))
    \mid z
  \right].
  \label{eq:pricing-residual}
\end{align}
The residual \(\mathcal P_\theta\) is zero when the lender breaks even.
For fixed current \(z\) and chosen \((k',b')\), the expectation integrates over
future shocks \(z'\), because those shocks determine whether the future state
\((k',b',z')\) is a default state or a solvent state.

\subsection{Policy maximization residual}

To keep the lecture self-contained, we use a direct maximization residual rather
than first-order conditions. The idea is simple: the chosen action should do at
least as well as any alternative action.

For any candidate alternative \(a=(\bar k',\bar b')\), define
\[
  \widehat Q_\theta(k,b,z;a)
\]
by taking the right side of \cref{eq:q-value}, replacing \(V\) and
\(\widetilde r\) by their neural-network approximations, and using the action
\(a\). The policy loss against the alternative \(a\) is
\begin{equation}
  \mathcal M_\theta(k,b,z;a)
  =
  \max\left\{
    0,\,
    \widehat Q_\theta(k,b,z;a)
    -
    \widehat Q_\theta(k,b,z;(\widehat k',\widehat b'))
  \right\}.
  \label{eq:policy-residual}
\end{equation}

\begin{proposition}[Direct policy residual]
\label{prop:direct-policy-loss}
Fix a state \((k,b,z)\). If
\[
  \mathcal M_\theta(k,b,z;a)=0
\]
for every feasible alternative action \(a=(\bar k',\bar b')\), then the
network-chosen action \((\widehat k',\widehat b')\) maximizes
\(\widehat Q_\theta(k,b,z;\cdot)\) over the feasible action set.
\end{proposition}

\begin{proof}
The equation \(\mathcal M_\theta(k,b,z;a)=0\) means
\[
  \widehat Q_\theta(k,b,z;a)
  -
  \widehat Q_\theta(k,b,z;(\widehat k',\widehat b'))
  \le 0
\]
for that alternative action \(a\). Therefore
\[
  \widehat Q_\theta(k,b,z;(\widehat k',\widehat b'))
  \ge
  \widehat Q_\theta(k,b,z;a).
\]
If this holds for every feasible \(a\), then no feasible alternative gives a
higher value. That is exactly the definition of a maximizer.
\end{proof}

\begin{remark}[Sampled alternatives are not global proof]
The proposition is an exact mathematical statement because it assumes every
feasible alternative action is checked. In computation, we usually test only a
finite sample of alternatives. Zero residuals on that finite test set certify
only that the chosen policy beats the sampled alternatives, not that it is a
global maximizer over all possible \((k',b')\).
\end{remark}

\subsection{Euler residuals from first-order conditions}
\label{subsec:euler-residuals}

Directly minimizing Bellman residuals can be numerically difficult. This is one
reason Euler-equation methods are common in dynamic economics and reinforcement
learning. Maliar et al. also describe Euler-residual minimization
\citep[Definitions 2.6 and 2.7]{maliar2021deep}. In the risky-debt model, Euler
residuals are useful, but only in the smooth part of the problem: the firm must
be continuing, the policy choice must be interior, and the risky rate and value
function must be differentiable. At default boundaries, borrowing limits, or
other kinks, one needs complementarity conditions, subgradients, or smoothing.

Let
\[
  m(e)=1+\eta'(e)
\]
be the marginal value of one more dollar of current cash flow. If external
finance is costly, then \(m(e)\) is high when the firm is short of cash. For an
interior continuation state, define
\[
  \bar e
  =
  e(k,k',b,b',z;\widetilde r(z,k',b')).
\]
The total derivative of current cash flow with respect to a choice must include
the effect of that choice on the risky rate:
\begin{align}
  \frac{d\bar e}{dk'}
  &=
  -(1+\psi_I(I,k))
  +e_{\widetilde r}\frac{\partial \widetilde r(z,k',b')}{\partial k'},
  \label{eq:cashflow-total-k}\\
  \frac{d\bar e}{db'}
  &=
  \frac{1}{1+\widetilde r}
  +\frac{\tau\widetilde r}{(1+\widetilde r)(1+r)}
  +e_{\widetilde r}\frac{\partial \widetilde r(z,k',b')}{\partial b'},
  \label{eq:cashflow-total-b}
\end{align}
where \(I=k'-(1-\delta)k\), \(\psi_I\) is the derivative of \(\psi\) with
respect to investment, and
\begin{equation}
  e_{\widetilde r}
  =
  \frac{b'}{(1+\widetilde r)^2}
  \left(-1+\frac{\tau}{1+r}\right).
  \label{eq:cashflow-rate-derivative}
\end{equation}
All terms in \cref{eq:cashflow-total-k,eq:cashflow-total-b,eq:cashflow-rate-derivative}
are evaluated at \((k,k',b,b',z)\).

\begin{proposition}[Interior first-order conditions]
\label{prop:interior-foc}
Suppose the current state is a continuation state, the optimal action
\((k',b')\) is interior, and the relevant functions are differentiable. Then the
optimal policy satisfies
\begin{align}
  0
  &=
  m(\bar e)\frac{d\bar e}{dk'}
  +\beta \E[V^\star_k(k',b',z')\mid z],
  \label{eq:foc-k}\\
  0
  &=
  m(\bar e)\frac{d\bar e}{db'}
  +\beta \E[V^\star_b(k',b',z')\mid z].
  \label{eq:foc-b}
\end{align}
\end{proposition}

\begin{proof}
At an interior continuation state, the firm chooses \((k',b')\) to maximize
\[
  e+\eta(e)+\beta \E[V^\star(k',b',z')\mid z].
\]
Differentiate this expression with respect to \(k'\) and \(b'\). The derivative
of \(e+\eta(e)\) is \((1+\eta'(e))de=m(e)de\). The derivative of the continuation
term is the expected derivative of \(V^\star\) with respect to the next state.
Setting both derivatives equal to zero gives
\cref{eq:foc-k,eq:foc-b}.
\end{proof}

The first-order conditions still contain derivatives of the value function.
Those derivatives can be obtained by automatic differentiation of a value
network. Alternatively, in smooth continuation states, we can use envelope
conditions to eliminate them.

\begin{proposition}[Envelope derivatives]
\label{prop:envelope-derivatives}
Away from the default boundary, write
\[
  h^{k\star}=h^{k\star}(k,b,z),
  \qquad
  h^{b\star}=h^{b\star}(k,b,z),
\]
and define
\[
  A(k,b,z)
  =
  (1-\tau)\pi_k(k,z)
  +(1-\delta)(1+\psi_I(I,k))
  -\psi_k(I,k),
\]
where \(I=h^{k\star}(k,b,z)-(1-\delta)k\). Then
\begin{align}
  V^\star_k(k,b,z)
  &=
  (1-D(k,b,z))\,m(e(k,h^{k\star},b,h^{b\star},z;\widetilde r))\,A(k,b,z),
  \label{eq:envelope-k}\\
  V^\star_b(k,b,z)
  &=
  -(1-D(k,b,z))\,m(e(k,h^{k\star},b,h^{b\star},z;\widetilde r)).
  \label{eq:envelope-b}
\end{align}
\end{proposition}

\begin{proof}
In a default state, \(V^\star=0\) locally, so the derivative is zero away from
the boundary. In a continuation state, \(V^\star=G^\star\). The envelope theorem
says that when differentiating with respect to the current state variables
\((k,b)\), we can hold the optimal choices fixed. Differentiating current cash
flow with respect to \(k\) gives \(A(k,b,z)\), and differentiating with respect
to beginning-of-period debt \(b\) gives \(-1\). Multiplying by the marginal value
of cash flow \(m(e)\) gives \cref{eq:envelope-k,eq:envelope-b}.
\end{proof}

Combining the first-order conditions and the envelope derivatives gives an
Euler-residual objective. For a network policy, let
\begin{equation}
  \widehat{\mathcal E}^k_\theta(k,b,z)
  =
  \widehat m_t\frac{d\widehat e_t}{dk'}
  +\beta \E[
    (1-\widehat D_{\theta,t+1})\widehat m_{t+1}\widehat A_{t+1}
    \mid z],
  \label{eq:euler-residual-k}
\end{equation}
and
\begin{equation}
  \widehat{\mathcal E}^b_\theta(k,b,z)
  =
  \widehat m_t\frac{d\widehat e_t}{db'}
  -\beta \E[
    (1-\widehat D_{\theta,t+1})\widehat m_{t+1}
    \mid z].
  \label{eq:euler-residual-b}
\end{equation}
Here \(\widehat e_t\) is current cash flow under the network policy,
\(\widehat m_t=1+\eta'(\widehat e_t)\), and the \(t+1\) quantities are computed
at the future state \((\widehat k',\widehat b',z')\) using the network's next
policy choices. The derivative terms \(d\widehat e_t/dk'\) and
\(d\widehat e_t/db'\) include the derivative of the risky rate network with
respect to the action. In practice, automatic differentiation can compute these
terms.

An Euler-style loss can replace the direct policy maximization loss on states
where the policy is smooth and interior:
\begin{align}
  \mathcal L^{\mathrm{Euler}}(\theta)
  =
  \E_{(k,b,z)\sim\mu}
  \Big[
    w_E^k(\widehat{\mathcal E}^k_\theta)^2
    +w_E^b(\widehat{\mathcal E}^b_\theta)^2
    +w_P\mathcal P_\theta^2
    +w_B\mathcal B_\theta^2
  \Big].
  \label{eq:euler-training-objective}
\end{align}
The Bellman/default residual \(\mathcal B_\theta\) is still useful because Euler
equations do not by themselves locate the default boundary or pin down the
level of equity value. In applications one may put most of the weight on the
Euler and pricing residuals in smooth continuation regions, and use
\(\mathcal B_\theta\) mainly to discipline limited liability and default. Thus
\cref{eq:euler-training-objective} should be read as the interior part of the
method, not as the full risky-debt training objective.

\subsection{Training with Euler residuals and nonsmooth regions}
\label{subsec:euler-kkt-training}

The practical solver should not train Euler residuals everywhere. Instead, it
should use different residuals in different regions of the state space. Let
\[
  C_\theta(k,b,z)=\1\{\widehat G_\theta(k,b,z)>0\}
\]
be the network's continuation indicator, and let
\[
  I_\theta(k,b,z)
\]
be an interior indicator that equals one when the policy choices are safely away
from imposed numerical bounds, such as \(\underline k<\widehat k'<\overline k\)
and \(\underline b<\widehat b'<\overline b\). The Euler residual should receive
full weight only when
\[
  S_\theta(k,b,z)=C_\theta(k,b,z)I_\theta(k,b,z)=1.
\]
In code, the sharp indicators above are often replaced by smooth gates. For
example, a smooth continuation gate can be
\[
  C_{\theta,\kappa}(k,b,z)
  =
  \frac{1}{1+\exp(-\kappa \widehat G_\theta(k,b,z))},
\]
where large \(\kappa\) makes the gate close to the sharp continuation
indicator.

If the firm continues but the action is not interior, the first-order condition
\(0=\partial Q/\partial a\) should be replaced by a
Karush-Kuhn-Tucker condition. If the firm defaults, the action FOC is not the
right condition, because shareholders are choosing the outside option with value
zero. That region is disciplined by the limited-liability Bellman residual.

For a scalar choice \(x\in[\underline x,\overline x]\), where \(x\) stands for
either \(k'\) or \(b'\), let
\[
  g_x(k,b,z)
  =
  \left.
  \frac{\partial \widehat Q_\theta(k,b,z;k',b')}
  {\partial x}
  \right|_{(k',b')=(\widehat k',\widehat b')}
\]
be the network-implied derivative with respect to that choice, holding the other
choice fixed at its network value. Here \(g_x\) is the derivative of the
objective being maximized, not the derivative of the negative objective. At a
maximum, the derivative must satisfy
\[
  g_x\le 0 \quad\text{at }x=\underline x,\qquad
  g_x=0 \quad\text{inside the interval},\qquad
  g_x\ge 0 \quad\text{at }x=\overline x.
\]
Equivalently, with lower and upper multipliers \(\lambda_x^L,\lambda_x^U\ge0\),
\begin{align}
  0&=g_x+\lambda_x^L-\lambda_x^U, \label{eq:kkt-stationarity}\\
  0&=\lambda_x^L(x-\underline x), \label{eq:kkt-lower}\\
  0&=\lambda_x^U(\overline x-x). \label{eq:kkt-upper}
\end{align}
These equations can be implemented by adding multiplier networks and using
Fischer-Burmeister residuals for the lower and upper bounds:
\begin{align}
  \mathcal K^x_\theta(k,b,z)
  &=
  \bigl(g_x+\lambda_x^L-\lambda_x^U\bigr)^2 \notag\\
  &\quad
  +\operatorname{FB}\bigl(\lambda_x^L,x-\underline x\bigr)^2
  +\operatorname{FB}\bigl(\lambda_x^U,\overline x-x\bigr)^2.
  \label{eq:kkt-residual-x}
\end{align}
The full action-bound residual is
\[
  \mathcal K_\theta=\mathcal K^{k'}_\theta+\mathcal K^{b'}_\theta,
\]
with each term evaluated at the network action. This is the same
complementarity idea used earlier for limited liability.

The training idea is now simple. At a safe interior continuation state, a good
policy should make the Euler residuals small. At a continuation state on an
action boundary, a good policy should make the KKT residual small instead. At a
default state, the policy action is economically irrelevant to shareholders
because they choose the outside option, so the Bellman/default residual is the
main condition. The debt-pricing residual is separate: whenever the chosen debt
is positive, lenders must break even under the default probabilities implied by
the same value and policy networks.

The hybrid Euler-KKT loss is therefore
\begin{align}
  \mathcal L^{\mathrm{hybrid}}(\theta)
  =
  \E_{(k,b,z)\sim\mu}
  \Big[
    &w_E S_\theta
      \big((\widehat{\mathcal E}^k_\theta)^2
      +(\widehat{\mathcal E}^b_\theta)^2\big) \notag\\
    &+w_K C_\theta(1-I_\theta)\mathcal K_\theta(k,b,z)
      +w_P\mathcal P_\theta(z,\widehat k',\widehat b')^2
      +w_B\mathcal B_\theta(k,b,z)^2
  \Big],
  \label{eq:hybrid-euler-kkt-objective}
\end{align}
where the KKT term is weighted only in continuation states. More explicitly, if
\[
  J_\theta(k,b,z)=\1\{\widehat b'_\theta(k,b,z)>0\}
\]
is a positive-debt gate, then many implementations use
\[
  w_PJ_\theta\mathcal P_\theta(z,\widehat k',\widehat b')^2
\]
in place of the ungated pricing term. If \(b'\le0\), the pricing term is either
omitted or replaced by the chosen cash-return convention; for positive debt
choices it is always enforced.
If the implementation uses action transformations that always keep
\(\widehat k'\) and \(\widehat b'\) strictly inside their bounds, then
\(\mathcal K_\theta\) may be unnecessary, but the Euler residual can become weak
near the artificial boundary because the transformation saturates.

\begin{proposition}[What the hybrid residuals certify]
\label{prop:hybrid-residuals}
Suppose the Bellman/default residual, pricing residual, Euler residuals on
smooth interior continuation states, and KKT residuals on constrained states are
all zero on a region. Then the approximating functions satisfy the local
necessary optimality conditions, limited liability, and debt-pricing conditions
on that region.
\end{proposition}

\begin{proof}
On smooth interior continuation states, zero Euler residuals are exactly the
first-order conditions from \cref{prop:interior-foc}. On constrained states,
zero KKT residuals are the first-order necessary conditions with active bounds.
Zero Bellman/default residual enforces the limited-liability default condition,
and zero pricing residual enforces lender break-even. Together these conditions
are the local necessary conditions for the risky-debt problem.
\end{proof}

\begin{remark}[Euler residuals are not a global optimality certificate]
\label{rem:euler-not-global}
Euler and KKT residuals are local conditions. They can miss a better distant
action when the objective is nonconcave. For that reason, a practical solver
should still test the trained policy against random alternative actions, report
out-of-sample Bellman and pricing residuals, and simulate the model to check
that default and leverage behavior are economically sensible.
\end{remark}

\subsection{Hard constraints as output transformations}
\label{subsec:hard-constraints}

The KKT residuals above are one way to handle action bounds. A different idea,
emphasized by \citet{hallhoffarth2023hard}, is to give constraints
\emph{lexicographic priority}: first restrict the neural network so that its
outputs are feasible by construction, and only then train the remaining
optimality residuals. In the language of the paper, these are
\emph{hard constraints}, in contrast to \emph{soft constraints} that are enforced
only through penalty terms such as Fischer-Burmeister losses.

The simplest examples are scalar bounds. If the raw neural network output is an
unrestricted number \(n_k(k,b,z)\), define
\begin{equation}
  \widehat k'
  =
  \underline k+(\overline k-\underline k)
  \sigma(n_k(k,b,z)),
  \qquad
  \sigma(x)=\frac{1}{1+\exp(-x)}.
  \label{eq:hard-k-bound}
\end{equation}
Then \(\underline k<\widehat k'<\overline k\) automatically. Similarly,
\begin{equation}
  \widehat b'
  =
  \underline b+(\overline b-\underline b)
  \sigma(n_b(k,b,z))
  \label{eq:hard-b-bound}
\end{equation}
automatically satisfies \(\underline b<\widehat b'<\overline b\). If a variable
only needs to be positive, a softplus transformation can be used:
\begin{equation}
  x=\epsilon+\log(1+\exp(n_x)),
  \qquad \epsilon>0.
  \label{eq:softplus-positive}
\end{equation}
For example, one may parameterize the risky gross rate as
\[
  1+\widehat r_\theta(z,k',b')
  =
  \epsilon_r+\log(1+\exp(n_r(z,k',b'))),
\]
which ensures a positive gross promised return.

The smooth sigmoid transformation is best interpreted as a way to impose
artificial numerical bounds. If a bound is a genuine economic constraint that
may bind in equilibrium, a strictly interior sigmoid may be too restrictive
because it cannot place the choice exactly on the boundary. In that case one
can use a transformation that allows binding, such as a clipped positive output,
or keep the KKT/complementarity residual for that economic constraint. Clipping
can create nondifferentiable or flat-gradient regions, so it is a practical
trade-off rather than a free improvement. This is the same reason
\citet{hallhoffarth2023hard} distinguishes hard feasibility from ordinary
penalty-based constraint handling.

\begin{proposition}[Hard numerical bounds remove numerical-bound KKT residuals]
\label{prop:hard-bounds-remove-kkt}
Suppose the policy networks are parameterized by
\cref{eq:hard-k-bound,eq:hard-b-bound}. Then the numerical bounds on \(k'\) and
\(b'\) are satisfied at every state and every parameter value. If the
transformation stays strictly inside artificial numerical bounds that are not
intended to bind as economic constraints, the usual interior Euler conditions
replace the numerical-bound KKT residuals on those controls.
\end{proposition}

\begin{proof}
The sigmoid function always lies between zero and one. Therefore
\(\underline k+(\overline k-\underline k)\sigma(n_k)\) always lies strictly
between \(\underline k\) and \(\overline k\), and the same argument applies to
\(b'\). Since the transformed choices are interior points of the numerical
interval, no lower or upper numerical bound is active. Thus, for those artificial
box bounds, there is no active bound multiplier to learn. The relevant local
optimality condition is the Euler first-order condition with respect to the
free raw network output, equivalently the chain-rule version of the interior
condition with respect to the transformed control when the transformation has a
nonzero derivative. Numerically, a saturating sigmoid can still make gradients
very small near the artificial boundary.
\end{proof}

This hard-constraint idea is useful whenever part of a model can be solved by
algebra. For example, if an accounting identity gives one variable as an exact
function of the others, it is usually better to build that identity into the
network output than to ask the network to rediscover it from a penalty term.
The principle is: do not ask the network to learn equations that can be
satisfied exactly by construction.

For the risky-debt model, hard constraints are useful but not sufficient. They
can enforce:
\begin{itemize}[leftmargin=2em]
  \item numerical bounds on \(k'\) and \(b'\);
  \item positivity of quantities such as gross risky returns;
  \item accounting identities or algebraic definitions that can be solved
  explicitly from other variables.
\end{itemize}
They cannot, by themselves, solve the default free-boundary condition
\[
  V(k,b,z)=\max\{0,G(k,b,z)\}.
\]
The reason is that default is not merely a bound on a single output. It is a
regime choice: in continuation states equity value equals going-concern value,
while in default states equity value equals zero. A transformation such as
\(\widehat V_\theta=\operatorname{softplus}(n_V)\) guarantees
\(\widehat V_\theta\ge0\), but it does not guarantee
\(\widehat V_\theta=\widehat G_\theta\) in continuation states or
\(\widehat V_\theta=0\) in default states. Therefore the Bellman/default
residual \(\mathcal B_\theta\) remains necessary. Even if one hard-codes
\(\widehat V_\theta=\max\{0,\widehat G_\theta\}\) with a smooth approximation,
one still has to learn the correct continuation object \(\widehat G_\theta\) and
the default boundary consistently with debt pricing.

The recommended practical objective is therefore a hard-constrained hybrid
loss:
\begin{align}
  \mathcal L^{\mathrm{hard\text{-}hybrid}}(\theta)
  =
  \E_{(k,b,z)\sim\mu}
  \Big[
    &w_E C_\theta
      \big((\widehat{\mathcal E}^k_\theta)^2
      +(\widehat{\mathcal E}^b_\theta)^2\big) \notag\\
    &+w_PJ_\theta
      \mathcal P_\theta(z,\widehat k',\widehat b')^2
      +w_B\mathcal B_\theta(k,b,z)^2
  \Big],
  \label{eq:hard-hybrid-objective}
\end{align}
where the action-bound KKT term is omitted because the relevant bounds are
already enforced by the output layer. If some genuine economic inequality is not
hard-coded, or if the implementation uses clipping that creates flat-gradient
regions, then the corresponding KKT or complementarity residual should be kept.

\begin{remark}[Hard constraints do not prove optimality]
\label{rem:hard-constraints-not-optimality}
Hard constraints guarantee feasibility, not optimality. They make it impossible
for the network to violate selected constraints, which can greatly improve
training stability. The Euler, pricing, Bellman/default, and validation
diagnostics are still needed to check whether the feasible policy is close to an
economic solution.
\end{remark}

\section{The training objective}

Let \(\mu\) be a distribution over states \((k,b,z)\). It may be a fixed training
distribution over a chosen rectangle, or it may be the simulated ergodic
distribution generated by the current policy, as in the stochastic simulation
logic emphasized by \citet[Section 3]{maliar2021deep}.

There are now two natural training choices. The first uses sampled alternative
actions to test optimality directly. Let \(\nu(\cdot\mid k,b,z)\) be a
distribution over alternative actions used to test whether the policy is close
to optimal. The theoretical risk is
\begin{align}
  \mathcal L(\theta)
  &=
  \E_{(k,b,z)\sim\mu}
  \Big[
    w_B \mathcal B_\theta(k,b,z)^2
    +w_P \mathcal P_\theta(z,\widehat k',\widehat b')^2
  \notag\\
  &\qquad\qquad
    +w_M
    \E_{a\sim\nu(\cdot\mid k,b,z)}
    [\mathcal M_\theta(k,b,z;a)^2]
  \Big],
  \label{eq:training-objective}
\end{align}
where \(w_B,w_P,w_M>0\) are weights chosen so that the different residuals have
comparable numerical scale. This is called a \emph{weighted scalarization}: many
requirements are collapsed into one number before the optimizer sees them. It is
simple, but the choice of weights can be difficult. The pricing residual,
Bellman/default residual, Euler residuals, and direct maximization residuals can
have different units, different numerical sizes, and different training speeds.

The second choice uses the hybrid Euler-KKT objective in
\cref{eq:hybrid-euler-kkt-objective}. This is often cheaper in high-dimensional
action spaces because it avoids searching over many alternative actions.
However, it gives local first-order evidence, not global maximization evidence.
If the policy outputs are hard-constrained as in
\cref{subsec:hard-constraints}, the hard-constrained hybrid objective in
\cref{eq:hard-hybrid-objective} can be used instead.

The empirical training loss with \(N\) sampled states for the direct
maximization version is
\begin{equation}
  \mathcal L_N(\theta)
  =
  \frac{1}{N}\sum_{i=1}^N
  \left[
    w_B \mathcal B_\theta(s_i)^2
    +w_P \mathcal P_\theta(z_i,\widehat k'_i,\widehat b'_i)^2
    +w_M \mathcal M_\theta(s_i;a_i)^2
  \right].
  \label{eq:empirical-objective}
\end{equation}
The empirical hybrid Euler-KKT loss replaces the policy-maximization term by
Euler and KKT terms:
\begin{align}
  \mathcal L^{\mathrm{hybrid}}_N(\theta)
  &=
  \frac{1}{N}\sum_{i=1}^N
  \Big[
    w_E S_\theta(s_i)
      \big((\widehat{\mathcal E}^k_\theta(s_i))^2
      +(\widehat{\mathcal E}^b_\theta(s_i))^2\big) \notag\\
  &\qquad
    +w_KC_\theta(s_i)(1-I_\theta(s_i))\mathcal K_\theta(s_i)
    +w_PJ_\theta(s_i)
      \mathcal P_\theta(z_i,\widehat k'_i,\widehat b'_i)^2 \notag\\
  &\qquad
    +w_B\mathcal B_\theta(s_i)^2
  \Big].
  \label{eq:empirical-hybrid-objective}
\end{align}
Training means approximately solving either
\begin{equation}
  \theta^\star\in \argmin_\theta \mathcal L_N(\theta).
  \label{eq:train-min}
\end{equation}
or the same minimization with \(\mathcal L^{\mathrm{hybrid}}_N\) in place of
\(\mathcal L_N\).
If hard output constraints are used, the empirical analogue of
\cref{eq:hard-hybrid-objective} is
\begin{align}
  \mathcal L^{\mathrm{hard\text{-}hybrid}}_N(\theta)
  &=
  \frac{1}{N}\sum_{i=1}^N
  \Big[
    w_E C_\theta(s_i)
      \big((\widehat{\mathcal E}^k_\theta(s_i))^2
      +(\widehat{\mathcal E}^b_\theta(s_i))^2\big) \notag\\
  &\qquad
    +w_PJ_\theta(s_i)
      \mathcal P_\theta(z_i,\widehat k'_i,\widehat b'_i)^2
    +w_B\mathcal B_\theta(s_i)^2
  \Big].
  \label{eq:empirical-hard-hybrid-objective}
\end{align}

\subsection{Replacing fixed weights by multi-objective training}
\label{subsec:multiobjective-training}

The fixed weights in
\cref{eq:training-objective,eq:empirical-hybrid-objective,eq:empirical-hard-hybrid-objective}
are often the least pleasant part of the computation. If \(w_P\) is too small,
the equity policy may look good while the debt price violates lender break-even.
If \(w_B\) is too small, the default boundary may drift. If \(w_E\) is too large,
the solver may chase local Euler equations while ignoring the level of the value
function. There is no universal scale that makes these choices automatic.

A practical alternative is \emph{multi-objective optimization} (MOO). Instead
of relying only on one permanent weighted sum, we keep the residual losses
separate and let the current gradients help decide the update direction. This
is the viewpoint of the multiple-gradient descent algorithm of
\citet{desideri2012mgda}; \citet{sener2018mtlmoo} use the same broad
multi-objective perspective for neural-network multi-task learning. For example,
with hard output constraints one may define the mini-batch objectives
\begin{align}
  \mathcal L_E(\theta)
  &=
  \frac{1}{N}\sum_{i=1}^N
  C_\theta(s_i)
  \left[
    \bigl(\widehat{\mathcal E}^k_\theta(s_i)\bigr)^2
    +\bigl(\widehat{\mathcal E}^b_\theta(s_i)\bigr)^2
  \right],
  \label{eq:moo-euler-loss}\\
  \mathcal L_P(\theta)
  &=
  \frac{1}{N}\sum_{i=1}^N
  J_\theta(s_i)
  \mathcal P_\theta(z_i,\widehat k'_i,\widehat b'_i)^2,
  \label{eq:moo-pricing-loss}\\
  \mathcal L_B(\theta)
  &=
  \frac{1}{N}\sum_{i=1}^N
  \mathcal B_\theta(s_i)^2.
  \label{eq:moo-bellman-loss}
\end{align}
If action-bound KKT terms, direct policy-check terms, or local-anchor terms are
needed, each can be added as another objective. Thus the optimizer sees a vector
\[
  \mathcal L(\theta)
  =
  \bigl(\mathcal L_1(\theta),\ldots,\mathcal L_m(\theta)\bigr),
\]
not just the scalar sum \(\sum_{j=1}^m w_j\mathcal L_j(\theta)\).
For gradient-based training, the gates \(C_\theta\) and \(J_\theta\) are usually
implemented as smooth approximations, or the derivative statement below is read
as applying away from the default and zero-debt kinks.

The intuition is simple. Each residual loss pulls the network parameters in a
direction. Fixed weights decide in advance which pull is louder. MOO computes
the pulls on the current mini-batch and combines them adaptively.

\begin{definition}[Common descent and Pareto stationarity]
\label{def:common-descent}
Let \(g_j(\theta)=\nabla_\theta \mathcal L_j(\theta)\) be the gradient of the
\(j\)-th objective. A direction \(p\) is a \emph{common descent direction} if
\[
  g_j(\theta)^\top p < 0
  \qquad\text{for every }j.
\]
This means that a small move from \(\theta\) in direction \(p\) lowers every
objective to first order. A point is \emph{first-order Pareto stationary} if no
common descent direction exists.
\end{definition}

The easiest MOO rule to write down is the minimum-norm gradient rule used in
MGDA \citep{desideri2012mgda}. Let
\[
  \Delta_m
  =
  \left\{
    \alpha\in\R^m:
    \alpha_j\ge0,\quad
    \sum_{j=1}^m\alpha_j=1
  \right\}
\]
be the probability simplex. At a training step, compute all objective gradients
\[
  g_1,\ldots,g_m.
\]
MGDA chooses
\begin{equation}
  \alpha^\star
  \in
  \argmin_{\alpha\in\Delta_m}
  \left\|
    \sum_{j=1}^m \alpha_j g_j
  \right\|^2,
  \qquad
  d^\star=\sum_{j=1}^m\alpha_j^\star g_j.
  \label{eq:mgda-alpha}
\end{equation}
Then the neural-network optimizer uses \(d^\star\) as the combined gradient, so
the parameter update is roughly
\[
  \theta_{\mathrm{new}}
  =
  \theta-\gamma d^\star
\]
for a learning rate \(\gamma>0\), possibly after Adam or another stochastic
optimizer rescales the coordinates.

\begin{proposition}[MGDA gives a balanced first-order update]
\label{prop:mgda-common-descent}
Suppose the objectives are differentiable at the current parameter vector
\(\theta\), and let \(d^\star\) be defined by \cref{eq:mgda-alpha}.
If \(d^\star=0\), then \(\theta\) is first-order Pareto stationary. If
\(d^\star\ne0\), then \(-d^\star\) is a common descent direction:
\begin{equation}
  g_j^\top(-d^\star)\le -\|d^\star\|^2<0
  \qquad\text{for every }j.
  \label{eq:mgda-common-descent}
\end{equation}
\end{proposition}

\begin{proof}
The vector \(d^\star\) is the point in the convex hull of
\(\{g_1,\ldots,g_m\}\) that is closest to zero. Take any other point \(y\) in
this convex hull. For \(0<t<1\), the point
\[
  d_t=(1-t)d^\star+ty
\]
is also in the convex hull. Since \(d^\star\) is the closest point to zero,
\(\|d_t\|^2\ge \|d^\star\|^2\). Expanding this inequality gives
\[
  2t(d^\star)^\top(y-d^\star)+t^2\|y-d^\star\|^2\ge0.
\]
Divide by \(t\) and let \(t\) become very small. We get
\[
  (d^\star)^\top y\ge \|d^\star\|^2.
\]
Each gradient \(g_j\) is itself a point in the convex hull, so
\[
  g_j^\top d^\star\ge \|d^\star\|^2.
\]
If \(d^\star\ne0\), multiplying by \(-1\) gives
\cref{eq:mgda-common-descent}.

Now suppose \(d^\star=0\). Then zero is a convex combination of the gradients:
\[
  0=\sum_{j=1}^m\alpha_j^\star g_j.
\]
If a common descent direction \(p\) existed, then \(g_j^\top p<0\) for every
\(j\), and therefore
\[
  0
  =
  \left(\sum_{j=1}^m\alpha_j^\star g_j\right)^\top p
  =
  \sum_{j=1}^m\alpha_j^\star g_j^\top p
  <0,
\]
which is impossible. Hence no common descent direction exists.
\end{proof}

This proposition is a first-order statement. It says what happens for an
infinitesimally small step when the objectives are smooth. It does not say that
a finite Adam step always lowers every loss on every mini-batch, and it does not
prove that a neural network has solved the economic model. It does, however,
explain why MOO is attractive: it reduces reliance on hand-tuned fixed weights
by using an adaptive combination based on the current gradient conflict.

A practical MGDA training step is as follows. First form one loss for each
economic condition. Then compute one gradient vector for each loss and build the
Gram matrix
\[
  G_{ij}=g_i^\top g_j,
\]
solve \cref{eq:mgda-alpha} for the simplex weights \(\alpha^\star\), and update
the network using the combined gradient \(d^\star\). The Gram matrix is useful
because it lets us compute the MGDA objective
\[
  \left\|\sum_{j=1}^m\alpha_j g_j\right\|^2
  =
  \alpha^\top G\alpha
\]
without writing out every coordinate of every gradient in the lecture note.
During training one should record diagnostics such as each objective loss,
heldout objective losses, gradient norms, the gradient cosine matrix, the
simplex weights \(\alpha_j^\star\), and the directional derivatives
\(g_j^\top d^\star\). These diagnostics matter because MOO reduces manual
weight tuning; it does not remove the need to monitor whether every economic
equation is improving. In practice, MOO should be read as reducing sensitivity
to weights, not abolishing all scale and normalization choices.

For the risky-debt solver, the practical recipe is:
\begin{enumerate}[leftmargin=2em]
  \item Use hard output transformations for simple feasibility requirements
  whenever possible.
  \item Split the remaining conditions into separate objectives: Euler
  residuals, debt-pricing residuals, Bellman/default residuals, KKT or
  complementarity residuals when needed, and direct policy-check or local-anchor
  residuals when global policy errors are a concern.
  \item Choose economically meaningful scales for the objective components.
  MOO reduces the need for fixed weights, but it does not magically fix bad
  units. For example, an Euler residual measured in marginal values and a
  pricing residual measured in promised debt payments should not be compared
  blindly without considering their natural economic scales.
  \item On each mini-batch, compute one gradient \(g_j\) for each objective and
  aggregate the gradients by the MGDA rule in \cref{eq:mgda-alpha}.
  \item Keep per-objective validation diagnostics. A small aggregate update is
  not enough; pricing, default, Euler, and action-check residuals should all be
  reported separately.
  \item Use economic veto rules when selecting a final network. For example, if
  the pricing residual or Bellman/default residual deteriorates beyond an
  acceptable tolerance on heldout states, reject the network even if an average
  scalar loss looks good.
\end{enumerate}

This makes the method more practical because the researcher no longer has to
guess a single permanent set of weights for all residuals. The remaining choices
are more interpretable: how to group the objectives, how to normalize them, what
heldout diagnostics must pass, and which MOO aggregation rule is stable for the
problem size.

\begin{proposition}[Zero direct residuals imply a solved approximate model]
\label{prop:zero-loss-solution}
Suppose the residuals
\[
  \mathcal B_\theta,\qquad \mathcal P_\theta,\qquad \mathcal M_\theta
\]
are all zero for every state in the training region and for every feasible
alternative action used in the maximization test. Then the approximating
functions \(\widehat V_\theta\), \(\widehat h^k_\theta\),
\(\widehat h^b_\theta\), and \(\widehat r_\theta\) satisfy the risky-debt
Bellman equation, the lender zero-profit condition, and the policy optimality
condition on that region.
\end{proposition}

\begin{proof}
If \(\mathcal B_\theta=0\), then by the Fischer-Burmeister equivalence the
approximated value function satisfies
\[
  \widehat V_\theta(k,b,z)
  =
  \max\{0,\widehat G_\theta(k,b,z)\}.
\]
This is the limited-liability Bellman equation for the approximating functions.

If \(\mathcal P_\theta=0\), then the expected payoff to the lender equals
\(b'(1+r)\). This is exactly the zero-profit debt-pricing condition.

If \(\mathcal M_\theta=0\) for every feasible alternative action, then
\cref{prop:direct-policy-loss} says that the network policy maximizes the
network-implied going-concern value. Combining the three statements gives the
claim.
\end{proof}

\section{Comparison with Duarte and Fonseca's structural-estimation algorithm}
\label{sec:duarte-fonseca-comparison}

A recent paper by \citet{duartefonseca2026ai} applies a closely related neural
idea to structural estimation. Their application is a dynamic corporate finance
model with leverage, investment, cash, endogenous default, costly equity
issuance, and non-convex adjustment costs. The paper is especially useful for
us because it studies a model with many of the same computational difficulties
as the risky-debt problem in this note: a value function, policy functions,
default, risky debt prices, and nonsmooth payoffs. Since the paper is a recent
working paper, we use it here as a carefully cited algorithmic reference, not
as a theorem that all such neural estimators must converge.

\subsection{What their algorithm does}

The algorithm in \citet[Section 3]{duartefonseca2026ai} has three main blocks.

First, it solves the firm's dynamic problem by neural policy iteration. The
value network is written as a function not only of the economic state, but also
of the structural parameters:
\[
  V(z,k,b,\beta;\Theta_V),
\]
where \(\beta\) denotes model parameters and \(\Theta_V\) denotes neural-network
weights. This is sometimes called a \emph{universal value function}: one network
represents many parameter-specific value functions. For a fixed policy, the
value-network weights are trained to reduce a squared Bellman residual
\citep[Section 3.1.1 and eqs. (14)--(15)]{duartefonseca2026ai}. Then, holding the
value network fixed, policy networks are trained to increase the right-hand side
of the Bellman equation \citep[Section 3.1.2 and eqs. (16)--(17)]{duartefonseca2026ai}.
The algorithm alternates between these two steps.

Second, the trained value and policy networks are used to simulate firm panels
for many parameter vectors. In their application, the authors refine the neural
solution on a grid before simulation. The refinement uses exact nonsmoothed
payoffs, local policy improvement, policy evaluation, and recomputation of bond
prices \citep[Section 3.2.1 and Appendix A.5]{duartefonseca2026ai}.

Third, the paper trains separate \emph{moment networks}. If \(m_j\) is the
\(j\)-th simulated moment, the moment network approximates
\[
  g_j(\beta) = \E[m_j\mid \beta].
\]
These networks turn simulated model moments into differentiable functions of
the parameters. Estimation then minimizes the usual weighted distance between
data moments and model moments,
\begin{equation}
  \widehat\beta
  \in
  \argmin_\beta
  \bigl(\widehat m-g(\beta)\bigr)^\top
  W
  \bigl(\widehat m-g(\beta)\bigr),
  \label{eq:duarte-moment-objective}
\end{equation}
using gradients of the moment networks
\citep[Section 3.3 and eqs. (18)--(21)]{duartefonseca2026ai}. For a general
equilibrium extension, they also treat the wage as a pseudo parameter and add a
market-clearing imbalance as another target moment
\citep[Section 3.4]{duartefonseca2026ai}. Finally, they use minimum-loss
functions to shrink parameter bounds and diagnose identification
\citep[Sections 3.5--3.6]{duartefonseca2026ai}.

\subsection{How this compares with the first approach in this note}

The first approach in this note is the direct value-and-policy approach: learn
\[
  \widehat V_\theta,\qquad
  \widehat h^k_\theta,\qquad
  \widehat h^b_\theta,\qquad
  \widehat r_\theta
\]
by minimizing Bellman/default, debt-pricing, and policy-maximization residuals.
At the level of the \emph{model solution block}, this is very close in spirit to
the first block of \citet{duartefonseca2026ai}. Both methods use neural networks
for value and policy functions. Both train with stochastic gradients and
automatic differentiation. Both can sample states instead of storing the value
function on a full grid. Both can smooth nonsmooth pieces during training and
then check accuracy with more direct policy or grid diagnostics.

So your intuition is essentially right, with one important qualification. If
our only goal is to solve one risky-debt model at one fixed parameter vector,
if we use the direct Bellman/policy objective rather than Euler residuals, and
if both methods impose the same lender zero-profit pricing equation, then the
solution step is not fundamentally different from the value/policy part of
\citet{duartefonseca2026ai}. The difference is mostly numerical organization:
this note writes the Bellman/default, pricing, and policy conditions as
residuals to be minimized jointly, while \citet{duartefonseca2026ai} present an
alternating policy-evaluation and policy-improvement procedure, with a lagged
bond-price target during training and a grid-refinement check afterward.

However, the full Duarte--Fonseca algorithm is more than a solver. It is a
structural-estimation pipeline. It solves across many parameter values at once,
simulates moments, trains moment-surrogate networks, estimates parameters using
\cref{eq:duarte-moment-objective}, treats general-equilibrium prices as pseudo
parameters when needed, and uses minimum-loss functions for global
identification diagnostics. Those parts are not contained in the basic
Maliar--Maliar--Winant solution method by themselves. They are built on top of
a neural solution method.

\begin{proposition}[Solver equivalence, but not estimation equivalence]
\label{prop:duarte-maliar-comparison}
For a fixed parameter vector and the same lender zero-profit pricing equation,
a direct Bellman/policy neural solver of the kind described in this note and
the value/policy block of \citet{duartefonseca2026ai} attack the same
mathematical fixed point: value, policy, default, and risky debt prices must be
mutually consistent. They differ mainly in training organization. But the full
Duarte--Fonseca structural-estimation algorithm contains additional moment,
estimation, market-clearing, and identification steps that are not implied by
solving the Bellman equation alone.
\end{proposition}

\begin{proof}
Both solver blocks try to find a value function and policy functions satisfying
the same type of Bellman optimality condition: the value equals current payoff
plus discounted expected continuation value under optimal controls, with
limited liability/default handled through a max or default rule. In a risky
debt model, the debt price must also satisfy the lender's zero-profit pricing
equation, because that price enters the firm's payoff and is itself determined
by future default. A direct residual method minimizes violations of these
conditions and violations of policy optimality. Neural policy iteration
alternates between fitting the value under the current policy and improving
the policy using the fitted value. If both procedures converge exactly and use
the same debt-pricing condition, they satisfy the same Bellman, pricing, and
policy optimality equations.

The moment-network and estimation steps are different objects. They take a
solved model, simulate model-implied moments, approximate the parameter-to-
moment map, and estimate parameters by matching moments. These steps are not
needed to define the firm's dynamic programming problem, and they do not follow
from solving the Bellman equation at one parameter vector. Hence the solver
blocks are closely related, while the full estimation algorithms are not the
same.
\end{proof}

\subsection{Is it suitable for the risky-debt problem?}

The method is promising for the risky-debt problem, especially if the researcher
wants to estimate parameters rather than merely compute one solution. The reason
is that risky-debt models are expensive to solve repeatedly. A neural solution
conditioned on parameters could amortize this cost: after training, one could
evaluate policies and simulate moments for many parameter values quickly.

For the narrower goal of solving the risky-debt model in this note, the method
is suitable only after addressing several gaps.

\begin{enumerate}[leftmargin=2em]
  \item \textbf{Debt-pricing fixed point.} In risky-debt models, the risky debt
  price depends on future default, while default depends on equity value, and
  equity value depends on the debt price. \citet[Appendix A.3.4]{duartefonseca2026ai}
  handle this circularity with a lagged target network for bond prices during
  training, followed by grid refinement. This is a useful numerical device, but
  for our note it should be treated as an approximation strategy, not a proof
  that the debt-pricing fixed point has been solved. We would still need
  out-of-sample pricing residuals.

  \item \textbf{Default boundary and nonsmooth payoffs.}
  \citet[Section 3.1.3 and Appendix A.4]{duartefonseca2026ai} smooth indicator
  functions during training and then refine using exact nonsmooth payoffs. This
  is natural for default and financing-cost kinks. The gap is that smoothing can
  bias the learned boundary, while grid refinement can become expensive as the
  number of state and control variables grows.

  \item \textbf{Local versus global policy improvement.} A neural policy
  improvement step can improve the policy where gradients point in the right
  direction, but the risky-debt objective can be nonconcave. A policy network
  may therefore settle near a local optimum. This note's direct policy-check
  residuals and the Duarte--Fonseca grid-refinement idea are both ways to test
  for better actions away from the network choice.

  \item \textbf{Parameter conditioning increases dimension.} Conditioning the
  value and policy networks on structural parameters is powerful for estimation,
  but it enlarges the input space. For the risky-debt model, parameters such as
  default costs, tax rates, adjustment-cost parameters, and financing-cost
  parameters may interact strongly. The network must be accurate not just at one
  parameter vector, but across the whole training region.

  \item \textbf{Moment networks add a second approximation.} If we use the full
  structural-estimation pipeline, then parameter estimates depend not only on
  value and policy approximation errors, but also on moment-network errors. The
  cross-validation and minimum-loss diagnostics in \citet[Sections 3.3 and
  3.6]{duartefonseca2026ai} are important, but they do not remove the need to
  validate estimates by direct simulation from the solved model.

  \item \textbf{General-equilibrium pseudo prices are not the same as risky debt
  price functions.} Treating a wage as a pseudo parameter works because the wage
  is a low-dimensional equilibrium object in their general equilibrium
  extension. In the risky-debt model, the debt price or risky rate is a function
  of state and promised actions, such as \(q(z,k',b')\) or
  \(\widetilde r(z,k',b')\). That object cannot usually be replaced by one
  scalar pseudo parameter. The lender zero-profit condition still needs to be
  enforced as a functional equation.
\end{enumerate}

Thus, for solving a fixed risky-debt model, Duarte--Fonseca should be viewed as
a close cousin of the direct Bellman/policy approach already in this note. For
estimating a risky-debt model, it suggests an important extension: condition the
solution on parameters, simulate moments, learn a differentiable moment map, and
then estimate parameters by matching moments. The main open work is to make the
default boundary, debt-pricing fixed point, nonconvex policy choice, and moment
surrogate errors visible in diagnostics rather than hidden inside a single
training loss.

\section{Parameter-conditioned Maliar--Maliar--Winant solvers}
\label{sec:parameter-conditioned-solver}

The preceding comparison was for one fixed parameter vector. But the neural
residual method of \citet{maliar2021deep} does not require the unknown
functions to depend only on economic states. We can enlarge the input of the
neural networks so that model parameters enter just like additional inputs.
This gives a \emph{parameter-conditioned} solver.

Let
\[
  \omega
  =
  (\delta,\tau,\alpha,\rho,\sigma_\varepsilon,\text{adjustment-cost parameters},
  \text{financing-cost parameters},\ldots)
\]
collect the structural parameters. Instead of learning one value function
\[
  V(k,b,z)
\]
for one fixed \(\omega\), we learn a family of value functions with one neural
network:
\[
  \widehat V_\theta(k,b,z,\omega).
\]
Likewise, we learn parameter-conditioned policy and pricing functions
\begin{equation}
  \widehat h^k_\theta(k,b,z,\omega),\qquad
  \widehat h^b_\theta(k,b,z,\omega),\qquad
  \widehat r_\theta(z,k',b',\omega).
  \label{eq:parameter-conditioned-networks}
\end{equation}
For each fixed value of \(\omega\), these functions should satisfy the same
Bellman, default, policy-optimality, and lender-pricing equations as before.
The training loss therefore becomes an average not only over states, but also
over parameters:
\begin{equation}
  \mathcal L(\theta)
  =
  \E_{(k,b,z,\omega)\sim \mu}
  \left[
    \mathcal B_\theta(k,b,z,\omega)^2
    +
    \mathcal P_\theta(k,b,z,\omega)^2
    +
    \mathcal M_\theta(k,b,z,\omega)^2
  \right],
  \label{eq:parameter-conditioned-loss}
\end{equation}
where \(\mu\) is the distribution from which we draw training points. Here
\(\mathcal B_\theta\) is the Bellman/default residual, \(\mathcal P_\theta\) is
the debt-pricing residual, and \(\mathcal M_\theta\) is the policy-improvement
or direct policy-check residual, now evaluated at the parameter vector
\(\omega\).

\begin{proposition}[Parameter conditioning gives the same solution block]
\label{prop:parameter-conditioned-equivalence}
Suppose a parameter-conditioned neural solver and the value/policy block of
\citet{duartefonseca2026ai} use the same state variables, the same structural
parameter vector, the same feasible action set, and the same lender
zero-profit pricing equation. Then both methods are trying to approximate the
same parameter-indexed fixed point
\[
  \omega
  \mapsto
  \bigl(V_\omega,h^k_\omega,h^b_\omega,\widetilde r_\omega\bigr).
\]
Thus, at the model-solution level, a parameter-conditioned
Maliar--Maliar--Winant residual solver is the same kind of object as the
Duarte--Fonseca neural solution block. The Duarte--Fonseca paper then adds
simulation, moment networks, estimation, and identification diagnostics on top
of this block.
\end{proposition}

\begin{proof}
Fix any parameter vector \(\omega\). Once \(\omega\) is fixed, the
parameter-conditioned networks in \cref{eq:parameter-conditioned-networks}
become ordinary functions of the economic state and promised actions. Setting
the residuals in \cref{eq:parameter-conditioned-loss} to zero at this fixed
\(\omega\) gives exactly the risky-debt Bellman/default equation, the
lender-pricing equation, and the policy-optimality condition for that
parameter vector. Duarte and Fonseca's value/policy block also seeks a value
function and policies indexed by parameters
\citep[Section 3.1 and eq. (13)]{duartefonseca2026ai}. Therefore both solve
the same family of parameter-specific dynamic programs. The remaining blocks
of their algorithm act after the model is solved: they simulate panels, learn
the parameter-to-moment map, and estimate parameters. Those blocks are not
part of the Bellman fixed point itself.
\end{proof}

This proposition is the clean answer to the question: yes, once the
Maliar--Maliar--Winant style solver is extended to take model parameters as
inputs, it becomes the same type of neural \emph{solution} object as the first
block of Duarte and Fonseca. The equality is not an equality of every numerical
detail. A residual solver may train all residuals jointly; neural policy
iteration may alternate policy evaluation and improvement. One may concatenate
states and parameters; another may use an architecture in which parameters
modulate the state network. These are different numerical designs for
approximating the same parameter-indexed solution map.

\subsection{Why parameter sampling is difficult}

The difficult new object in \cref{eq:parameter-conditioned-loss} is the
sampling distribution \(\mu\). For a fixed model, we only have to decide which
states \((k,b,z)\) to sample. For a parameter-conditioned model, we must decide
which pairs
\[
  (k,b,z,\omega)
\]
to sample. This is much harder.

\begin{enumerate}[leftmargin=2em]
  \item \textbf{The dimension grows quickly.} If the original risky-debt model
  has three state variables and we add eight parameters, the input space has
  eleven dimensions. A training set that looks dense in three dimensions can
  be extremely sparse in eleven dimensions. This is the usual curse of
  dimensionality, even though we are using neural networks rather than grids.
  Duarte and Fonseca explicitly emphasize this enlargement of the input space
  in their solution block \citep[Section 3.1]{duartefonseca2026ai}.

  \item \textbf{Parameters should not be sampled independently without thought.}
  A rectangular box such as
  \[
    \omega_j\in[\underline\omega_j,\overline\omega_j]
  \]
  is easy to sample from, but many combinations inside the box may be
  economically unimportant, weakly identified, or numerically pathological.
  For example, high debt benefits, low recovery, high productivity volatility,
  and low adjustment costs may create much more default activity than other
  combinations. Uniform sampling over a wide box can spend most computation on
  parameter regions that will never matter for estimation.

  \item \textbf{The relevant state region depends on parameters.} The typical
  range of capital, debt, and productivity is not the same for all
  \(\omega\). If shocks are very persistent, the relevant range of \(z\) may be
  wider. If adjustment costs are high, capital moves less. If default costs are
  low, high-debt states may be more common. Therefore the state-sampling rule
  should usually depend on the sampled parameter vector:
  \[
    \omega \sim \mu_\omega,
    \qquad
    (k,b,z)\sim \mu_x(\cdot\mid \omega).
  \]
  Duarte and Fonseca use parameter-dependent state bounds in their training
  procedure \citep[Appendix A.2 and Appendix A.3.1]{duartefonseca2026ai}.

  \item \textbf{Some parameters change the shape of the solution sharply.}
  Near a default boundary, a small change in recovery, volatility, or tax
  benefits can move the boundary a lot. A network trained on a coarse parameter
  sample may average across these changes and produce a smoothed default rule
  in the wrong place. This is especially dangerous for risky debt because the
  default boundary feeds back into the debt price.

  \item \textbf{Accuracy is not needed equally everywhere.} If the goal is
  estimation, the solver must be accurate near parameter values that match the
  data. It may be less important to solve remote parameter regions precisely.
  But we do not know the important region before estimating the model. This is
  the chicken-and-egg problem in parameter sampling. Duarte and Fonseca address
  this by starting from wide parameter bounds, sampling parameters inside those
  bounds, and then narrowing the bounds using minimum-loss functions with
  safeguards against excluding plausible or weakly identified regions
  \citep[Sections 3.5--3.6 and Appendix A.9]{duartefonseca2026ai}.

  \item \textbf{Weak identification makes shrinking dangerous.} Suppose a
  parameter barely affects the targeted moments. Then many values of that
  parameter fit the data almost equally well. If we shrink the training region
  too aggressively, we may accidentally throw away valid parameter values.
  Parameter sampling is therefore not only a numerical issue; it is also an
  identification issue.
\end{enumerate}

A practical parameter-conditioned solver should therefore report diagnostics
by parameter region. For example, after training, one should evaluate Bellman
residuals, pricing residuals, direct policy checks, and simulation moments on a
held-out sample of parameter vectors. It is not enough to report one average
loss over all sampled \((k,b,z,\omega)\). A small average can hide large errors
near default boundaries, high-leverage regions, or weakly identified parameter
directions.

\section{An alternative estimation route: state-space models, filtering, and HMC}
\label{sec:ssm-hmc-estimation}

Duarte and Fonseca's approach estimates parameters by learning a differentiable
map from parameters to simulated moments. There is another natural route. After
we solve the risky-debt model, the solved policy functions give laws of motion
for capital and debt. Together with the shock process, these laws of motion can
be treated as a \emph{state-space model}. Then we can estimate structural
parameters using filtering and Bayesian simulation methods such as Hamiltonian
Monte Carlo (HMC).

This is the same broad logic used in many likelihood-based dynamic models:
first solve the economic decision problem, then use the solved model as the
transition law for hidden economic states. State-space likelihood methods are a
standard way to handle latent states \citep{durbinkoopman2012ssm}. HMC is a
standard gradient-based method for sampling continuous posterior distributions
\citep{neal2011hmc}. When the state-space model is nonlinear or non-Gaussian,
particle filters and particle MCMC methods are often used
\citep{andrieu2010pmcmc}.

\subsection{From solved policies to a state-space model}

Let \(\omega\) denote the structural parameter vector. Suppose that, for each
\(\omega\), the neural solver gives approximate policy and pricing functions
\[
  \widehat h^k_\omega(k,b,z),\qquad
  \widehat h^b_\omega(k,b,z),\qquad
  \widehat r_\omega(z,k',b').
\]
For firm \(i\) at date \(t\), define the latent state
\[
  x_{i,t}=(k_{i,t},b_{i,t},z_{i,t},a_{i,t}),
\]
where \(a_{i,t}\in\{0,1\}\) is an indicator for whether the firm is active. If
the firm is active, the solved policies imply
\begin{align}
  k_{i,t+1}
  &=
  \widehat h^k_\omega(k_{i,t},b_{i,t},z_{i,t}), \label{eq:ssm-k-law}\\
  b_{i,t+1}
  &=
  \widehat h^b_\omega(k_{i,t},b_{i,t},z_{i,t}), \label{eq:ssm-b-law}\\
  \log z_{i,t+1}
  &=
  \rho \log z_{i,t}+\sigma_\varepsilon\varepsilon_{i,t+1},
  \qquad \varepsilon_{i,t+1}\sim N(0,1). \label{eq:ssm-z-law}
\end{align}
Default is determined by the solved continuation value. For example, using the
going-concern value \(\widehat G_\omega\) from \cref{def:going-concern-value},
one can write the next-period default indicator as
\begin{equation}
  d_{i,t+1}
  =
  \1\{\widehat G_\omega(k_{i,t+1},b_{i,t+1},z_{i,t+1})<0\}.
  \label{eq:ssm-default-indicator}
\end{equation}
If the firm defaults, the econometrician must specify what happens in the data:
the firm may exit the sample, be replaced by a new draw, or remain observed
with a default flag. This choice is not part of the Bellman equation; it is part
of the statistical observation model.

Let \(y_{i,t}\) denote observed data, such as book assets, leverage, investment,
equity issuance, market value, or a default indicator. A simple measurement
equation is
\begin{equation}
  y_{i,t}
  =
  m_\omega(x_{i,t})+u_{i,t},
  \qquad
  u_{i,t}\sim N(0,\Sigma_u),
  \label{eq:ssm-measurement}
\end{equation}
where \(m_\omega\) maps the model state into observable quantities. If the data
include a binary default observation, then that component is better written as
a Bernoulli measurement equation rather than a Gaussian one.

\begin{proposition}[Solved risky-debt model as a state-space model]
\label{prop:risky-debt-ssm}
Fix a parameter vector \(\omega\). If the policy functions, debt-pricing
function, shock law, default rule, and measurement equation are specified, then
the risky-debt model defines a state-space model
\[
  p_\omega(x_{t+1}\mid x_t),
  \qquad
  p_\omega(y_t\mid x_t).
\]
Therefore the likelihood of a time-series or panel dataset can be written as
\begin{equation}
  p_\omega(y_{1:T})
  =
  \int
  p_\omega(x_1)
  \prod_{t=1}^T
    p_\omega(y_t\mid x_t)
  \prod_{t=1}^{T-1}
    p_\omega(x_{t+1}\mid x_t)
  \,dx_{1:T}.
  \label{eq:ssm-likelihood}
\end{equation}
\end{proposition}

\begin{proof}
Equations \cref{eq:ssm-k-law,eq:ssm-b-law,eq:ssm-z-law} describe how the hidden
state evolves. The next state depends on the current state and the new shock,
so it defines a transition law \(p_\omega(x_{t+1}\mid x_t)\), allowing for
exit, replacement, or default flags if those are included in the state. Some
parts of this law may be deterministic: for example, conditional on the current
state and the shock, next period's capital and debt are pinned down by the
policy functions.
Equation \cref{eq:ssm-measurement} describes the distribution of observables
conditional on the hidden state, so it defines \(p_\omega(y_t\mid x_t)\). The
joint density of states and observations is therefore the initial density times
all transition densities times all measurement densities. Integrating out the
unobserved states gives \cref{eq:ssm-likelihood}.
\end{proof}

\subsection{Filtering and HMC}

The filtering recursion updates beliefs about the hidden state as data arrive:
\begin{align}
  p_\omega(x_t\mid y_{1:t-1})
  &=
  \int
  p_\omega(x_t\mid x_{t-1})
  p_\omega(x_{t-1}\mid y_{1:t-1})
  \,dx_{t-1}, \label{eq:ssm-predict}\\
  p_\omega(x_t\mid y_{1:t})
  &\propto
  p_\omega(y_t\mid x_t)
  p_\omega(x_t\mid y_{1:t-1}). \label{eq:ssm-update}
\end{align}
In a linear Gaussian model these equations reduce to the Kalman filter. The
risky-debt model is not linear Gaussian because policies are nonlinear and
default creates a discrete event. Therefore a particle filter is usually the
more natural filtering tool.

For Bayesian estimation, place a prior \(p(\omega)\) on the structural
parameters. The posterior is
\begin{equation}
  p(\omega\mid y_{1:T})
  \propto
  p(\omega)\,p_\omega(y_{1:T}).
  \label{eq:ssm-posterior}
\end{equation}
If \(p_\omega(y_{1:T})\) and its gradient can be evaluated accurately and
smoothly, HMC can sample from \cref{eq:ssm-posterior}. This is attractive
because HMC uses gradients to move efficiently through a high-dimensional
continuous parameter space.

\subsection{Can we do this for the risky-debt model?}

Yes, but it is harder than writing down \cref{eq:ssm-posterior}. The main
difficulties are the following.

\begin{enumerate}[leftmargin=2em]
  \item \textbf{The likelihood needs more than the solved Bellman equation.}
  Solving the firm problem gives policy functions. It does not by itself give
  measurement errors, replacement rules after default, sample-selection rules,
  or the distribution of initial firm states. These are statistical modeling
  choices.

  \item \textbf{Default makes the model non-Gaussian and partly discrete.}
  HMC works naturally for continuous differentiable posterior densities. A hard
  default indicator creates kinks and discrete events. One can either smooth
  the default rule for differentiable HMC, integrate default events with a
  filter, or use particle MCMC methods that are designed for nonlinear and
  non-Gaussian state-space models.

  \item \textbf{A particle-filter likelihood is noisy.} Particle filters give a
  simulation-based likelihood estimate. That estimate can be useful, but its
  Monte Carlo noise can make gradient-based HMC unstable unless the filter is
  carefully designed. Particle MCMC methods are a more direct fit when the
  likelihood is estimated by particles \citep{andrieu2010pmcmc}.

  \item \textbf{The neural solution depends on parameters.} In a likelihood or
  posterior evaluation at a new \(\omega\), the policies
  \(\widehat h^k_\omega\), \(\widehat h^b_\omega\), and prices
  \(\widehat r_\omega\) must be accurate at that \(\omega\). This again points
  toward a parameter-conditioned solver and held-out residual diagnostics by
  parameter region.

  \item \textbf{Approximation error enters the likelihood.} The likelihood in
  \cref{eq:ssm-likelihood} is the likelihood of the \emph{approximated} solved
  model. If the neural policy is wrong near the default boundary, the filter
  may fit the data by moving parameters in a biased direction. A serious
  estimation exercise should therefore report both filtering diagnostics and
  solver residual diagnostics.
\end{enumerate}

Thus the answer is yes in principle: after solving the risky-debt model, we can
turn it into a state-space model and estimate it with filtering plus HMC or
particle MCMC. This is different from Duarte and Fonseca's moment-network route.
Their route learns a parameter-to-moment map and matches moments. The
state-space route writes a likelihood for the time-series or panel data and
estimates from that likelihood. The price of the likelihood route is that we
must specify a full observation model and handle the nonlinear, nonsmooth, and
partly discrete default structure.

\section{A full-scale state-space specification}
\label{sec:full-scale-ssm}

This section gives one complete specification. It is not the only possible
specification, but it is complete enough that the likelihood is well defined.
It is designed to address the five problems listed above.

\subsection{Data and parameters}

Suppose we observe a panel of firms \(i=1,\ldots,N\) over dates
\(t=1,\ldots,T_i\). Let the observed vector be
\[
  y_{i,t}
  =
  \bigl(
    \log A_{i,t}^{obs},
    L_{i,t}^{obs},
    I_{i,t}^{obs}/A_{i,t}^{obs},
    Q_{i,t}^{obs},
    D_{i,t}^{obs}
  \bigr),
\]
where \(A^{obs}\) is an accounting measure of assets, \(L^{obs}\) is leverage,
\(I^{obs}/A^{obs}\) is the investment rate, \(Q^{obs}\) is a market-value or
valuation ratio, and \(D^{obs}\in\{0,1\}\) is an observed default or exit
indicator. If some components are missing, the likelihood simply omits those
components for that firm-date.

Let the estimated parameter vector be
\[
  \phi=(\omega,\Sigma_y,\mu_0,\Sigma_0,\pi_{01},\pi_{10},\kappa).
\]
Here \(\omega\) contains the economic parameters in the risky-debt model,
\(\Sigma_y\) contains measurement-error variances, \((\mu_0,\Sigma_0)\)
describe the initial distribution of latent states, \(\pi_{01}\) and
\(\pi_{10}\) are default-observation error probabilities, and \(\kappa>0\) is a
smoothing parameter used only in the differentiable approximation described
below. If one uses the exact discrete default model, \(\kappa\) can be omitted.

\subsection{Latent state and initial distribution}

For each firm, define the latent state
\[
  x_{i,t}=(k_{i,t},b_{i,t},z_{i,t},a_{i,t}),
\]
where \(a_{i,t}=1\) means the firm is active and \(a_{i,t}=0\) means the firm
has defaulted or exited. The initial state distribution is
\begin{align}
  \log k_{i,1} &\sim N(\mu_{k,0},\sigma_{k,0}^2),\\
  b_{i,1}/k_{i,1} &\sim N(\mu_{\ell,0},\sigma_{\ell,0}^2),\\
  \log z_{i,1} &\sim N(\mu_{z,0},\sigma_{z,0}^2),\\
  a_{i,1} &= 1.
\end{align}
These normal distributions can be replaced by truncated normal distributions
if the researcher wants to rule out economically impossible initial states.

\subsection{Solved policy transition}

For a proposed economic parameter vector \(\omega\), use the
parameter-conditioned neural solver to obtain
\[
  \widehat V_\omega,\qquad
  \widehat G_\omega,\qquad
  \widehat h^k_\omega,\qquad
  \widehat h^b_\omega,\qquad
  \widehat r_\omega.
\]
For an active firm, draw a productivity innovation
\(\varepsilon_{i,t+1}\sim N(0,1)\) and set
\begin{align}
  k^*_{i,t+1}
  &=
  \widehat h^k_\omega(k_{i,t},b_{i,t},z_{i,t}),\\
  b^*_{i,t+1}
  &=
  \widehat h^b_\omega(k_{i,t},b_{i,t},z_{i,t}),\\
  \log z^*_{i,t+1}
  &=
  \rho\log z_{i,t}+\sigma_\varepsilon\varepsilon_{i,t+1}.
\end{align}
The exact model default indicator is
\begin{equation}
  d^*_{i,t+1}
  =
  \1\{\widehat G_\omega(k^*_{i,t+1},b^*_{i,t+1},z^*_{i,t+1})<0\}.
  \label{eq:full-ssm-exact-default}
\end{equation}
If \(d^*_{i,t+1}=0\), then
\[
  x_{i,t+1}=(k^*_{i,t+1},b^*_{i,t+1},z^*_{i,t+1},1).
\]
If \(d^*_{i,t+1}=1\), use an absorbing exit rule:
\[
  a_{i,t+1}=0.
\]
All later observations for that firm are treated as missing except the default
or exit indicator. This absorbing-exit rule is natural for firm-level data in
which defaulted firms leave the sample. If instead the model is being simulated
as a stationary artificial panel, replace the defaulted firm by a new draw from
the initial distribution. The important point is that the rule must be stated
before estimation.

\subsection{Measurement equations}

For an active firm, define model-implied observables by
\begin{align}
  \log A_{i,t}^{mod}
  &=
  \log k_{i,t},\\
  L_{i,t}^{mod}
  &=
  b_{i,t}/k_{i,t},\\
  (I/A)_{i,t}^{mod}
  &=
  \frac{\widehat h^k_\omega(k_{i,t},b_{i,t},z_{i,t})
  -(1-\delta)k_{i,t}}{k_{i,t}},\\
  Q_{i,t}^{mod}
  &=
  \frac{\widehat V_\omega(k_{i,t},b_{i,t},z_{i,t})+b_{i,t}}{k_{i,t}}.
\end{align}
The last equation uses equity plus debt divided by capital as a simple
model-implied valuation ratio. If the data define \(Q\) differently, then this
measurement equation should be changed to match the data definition.

The continuous observations satisfy
\begin{equation}
  y^{cont}_{i,t}
  =
  m^{cont}_\omega(x_{i,t})
  +
  u_{i,t},
  \qquad
  u_{i,t}\sim N(0,\Sigma_y).
  \label{eq:full-ssm-cont-measurement}
\end{equation}
For the binary default observation, allow small classification errors:
\begin{equation}
\begin{aligned}
  \Pr(D_{i,t+1}^{obs}=1\mid d^*_{i,t+1}=0) &= \pi_{01},\\
  \Pr(D_{i,t+1}^{obs}=0\mid d^*_{i,t+1}=1) &= \pi_{10}.
\end{aligned}
\label{eq:full-ssm-default-measurement}
\end{equation}
These probabilities keep the likelihood from becoming exactly zero when data
and model default dates disagree by a small amount.

\subsection{Exact filtering specification}

The exact-discrete version uses \cref{eq:full-ssm-exact-default} as written. A
particle filter approximates
\[
  p_\phi(x_{i,t}\mid y_{i,1:t})
\]
by particles \(x_{i,t}^{(m)}\), \(m=1,\ldots,M\). At each date:
\begin{enumerate}[leftmargin=2em]
  \item Propagate each particle using the solved policy transition and a new
  shock draw.
  \item Compute the observation likelihood from
  \cref{eq:full-ssm-cont-measurement,eq:full-ssm-default-measurement}.
  \item Reweight particles by the observation likelihood.
  \item Resample if the effective sample size is too small.
\end{enumerate}
This gives an estimated likelihood
\[
  \widehat p_\phi(y_{i,1:T_i})
\]
for each firm. If firms are conditionally independent given \(\phi\), the panel
likelihood is
\begin{equation}
  \widehat p_\phi(Y)
  =
  \prod_{i=1}^N
  \widehat p_\phi(y_{i,1:T_i}).
  \label{eq:full-ssm-panel-likelihood}
\end{equation}

The Bayesian posterior is
\begin{equation}
  p(\phi\mid Y)
  \propto
  p(\phi)\,p_\phi(Y).
  \label{eq:full-ssm-panel-posterior}
\end{equation}
Because the likelihood is estimated by particles, the most faithful sampler is
a particle MCMC method, such as particle marginal Metropolis-Hastings
\citep{andrieu2010pmcmc}. This directly addresses the discreteness of default
and the non-Gaussian state-space likelihood. Plain HMC is not the safest
choice for this exact-discrete likelihood because resampling and hard default
events make the likelihood rough.

\subsection{Smoothed HMC specification}

If the researcher wants HMC, use a differentiable approximation. Replace the
hard default indicator by the smooth default probability
\begin{equation}
  p^D_{\kappa,i,t+1}
  =
  \frac{1}{1+\exp\{\widehat G_\omega(k^*_{i,t+1},b^*_{i,t+1},z^*_{i,t+1})/\kappa\}}.
  \label{eq:full-ssm-soft-default}
\end{equation}
When \(\widehat G_\omega\) is very negative, the default probability is close to
one. When \(\widehat G_\omega\) is very positive, it is close to zero. As
\(\kappa\downarrow0\), this approaches the hard default rule.

In the smoothed likelihood, use
\[
  D_{i,t+1}^{obs}
  \sim
  \operatorname{Bernoulli}(p^D_{\kappa,i,t+1})
\]
or combine \(p^D_{\kappa,i,t+1}\) with the classification-error probabilities
\(\pi_{01}\) and \(\pi_{10}\). If the resulting filtering likelihood is
differentiable in \(\phi\), HMC can target the approximate posterior
\begin{equation}
  \widetilde p_\kappa(\phi\mid Y)
  \propto
  p(\phi)\,\widetilde p_{\phi,\kappa}(Y).
  \label{eq:full-ssm-soft-posterior}
\end{equation}
The smoothing parameter is not harmless. A serious analysis should repeat the
estimation for smaller \(\kappa\) and check whether posterior conclusions are
stable. If conclusions change materially as \(\kappa\) changes, the smoothed
HMC posterior is a numerical approximation, not reliable evidence about the
exact default model.

\subsection{How this specification addresses the five problems}

\begin{enumerate}[leftmargin=2em]
  \item \textbf{More than the Bellman equation.} The model now specifies
  initial states, transitions, default/exit, measurement errors, default
  observation errors, and missing-data handling. The likelihood is therefore a
  statistical object, not merely a solved dynamic program.

  \item \textbf{Default discreteness.} The exact version uses a particle filter
  and particle MCMC, which are designed for nonlinear and non-Gaussian
  state-space models. The HMC version uses a smooth default probability and
  must be checked against the exact version.

  \item \textbf{Particle-filter noise.} The exact version uses the
  particle-filter likelihood inside particle MCMC rather than pretending that
  the likelihood is deterministic. In practice one should increase the number
  of particles until the log-likelihood estimate is stable enough for the
  sampler being used.

  \item \textbf{Parameter-dependent neural solution.} The solver must be
  parameter-conditioned on \(\omega\). During estimation, the posterior draws of
  \(\omega\) should remain inside the region where held-out Bellman, pricing,
  and policy-check residuals are small.

  \item \textbf{Approximation error.} Before accepting posterior draws, compute
  a solver diagnostic such as
  \[
    R_{0.95}(\omega)
    =
    \text{95th percentile of held-out solver residuals at }\omega.
  \]
  If \(R_{0.95}(\omega)\) is too large in posterior regions, the solution
  network must be retrained or refined there. The likelihood should not be
  trusted in parameter regions where the economic model has not actually been
  solved accurately.
\end{enumerate}

\begin{proposition}[A well-defined full-scale estimator]
\label{prop:full-scale-estimator}
Once the structural parameters, solved policy functions, initial-state law,
transition law, default/exit rule, measurement equations, and prior are
specified as above, the risky-debt model defines either an exact-discrete
particle-filter likelihood \(p_\phi(Y)\) or a smoothed likelihood
\(\widetilde p_{\phi,\kappa}(Y)\). Therefore estimation by particle MCMC or by
HMC on the smoothed posterior is mathematically well defined.
\end{proposition}

\begin{proof}
The initial-state law gives \(p_\phi(x_{i,1})\). The solved policy transition,
shock law, and default/exit rule give \(p_\phi(x_{i,t+1}\mid x_{i,t})\). The
continuous and binary measurement equations give \(p_\phi(y_{i,t}\mid x_{i,t})\).
Multiplying these terms and integrating over latent states gives a likelihood.
Combining the likelihood with a prior gives a posterior. Particle MCMC targets
the posterior associated with the exact-discrete state-space likelihood, while
HMC targets the posterior associated with the differentiable smoothed
likelihood.
\end{proof}

\section{Likelihood filtering versus moment matching}
\label{sec:likelihood-vs-moments}

We now have two estimation routes for the risky-debt model.

The \emph{moment-matching route} solves the model for many parameter values,
simulates firm panels, computes selected moments, learns a differentiable
moment map, and estimates parameters by matching data moments. This is the
route emphasized by \citet[Section 3]{duartefonseca2026ai}. In our notation,
it estimates
\begin{equation}
  \widehat\omega_{\mathrm{MM}}
  \in
  \argmin_\omega
  \bigl(\widehat m-g(\omega)\bigr)^\top W
  \bigl(\widehat m-g(\omega)\bigr),
  \label{eq:mm-estimator}
\end{equation}
where \(\widehat m\) are empirical moments and \(g(\omega)\) is the
model-implied moment function or its neural approximation.

The \emph{likelihood-filtering route} solves the model, embeds the solved
policies in a state-space model, filters latent states, and estimates from the
likelihood. In Bayesian form it estimates the posterior
\begin{equation}
  p(\phi\mid Y)
  \propto
  p(\phi)\,p_\phi(Y),
  \label{eq:likelihood-route-posterior}
\end{equation}
where \(p_\phi(Y)\) is obtained by a Kalman filter in special linear Gaussian
cases, by a particle filter in nonlinear non-Gaussian cases, or by a smoothed
differentiable approximation when using HMC.

\begin{proposition}[Different statistical objects]
\label{prop:mm-vs-likelihood-object}
Moment matching and likelihood filtering use the same solved economic model,
but they estimate different statistical objects. Moment matching compares a
small vector of selected features of the data with the same features in
simulated data. Likelihood filtering specifies a full probability law for the
observed panel and uses that law to evaluate the probability of the entire
observed path.
\end{proposition}

\begin{proof}
In \cref{eq:mm-estimator}, the data enter only through \(\widehat m\). If two
datasets have the same selected moments, the objective gives them the same
value. In \cref{eq:likelihood-route-posterior}, the data enter through the full
likelihood \(p_\phi(Y)\). The likelihood depends on the order of observations,
the timing of default, the path of latent states, and measurement errors. Thus
the two routes use the solved model differently: one through selected moments,
the other through a probability model for the full observed history.
\end{proof}

\subsection{Advantages of moment matching}

Moment matching is often the more forgiving route.

\begin{enumerate}[leftmargin=2em]
  \item \textbf{It does not require a complete observation model.} The
  researcher chooses moments such as mean leverage, investment volatility,
  default frequency, or autocorrelations. One does not need to specify every
  measurement-error distribution or the exact probability of every observed
  firm path.

  \item \textbf{It can target economically important features directly.} If
  leverage, default frequency, and investment moments are the main objects of
  interest, the estimator can put them directly in \(\widehat m\). This is
  useful for teaching and for empirical work where the model is intended to
  explain selected corporate-finance facts.

  \item \textbf{It can be computationally efficient after the surrogate is
  trained.} Once the moment network \(g(\omega)\) is accurate, evaluating
  \cref{eq:mm-estimator} is fast. This is the central computational advantage
  of the Duarte--Fonseca approach.

  \item \textbf{It is less exposed to small measurement-model mistakes.} A
  full likelihood can be very sensitive to a badly specified measurement error.
  Moment matching can be more robust if the selected moments are well chosen.

  \item \textbf{Identification diagnostics can be visual.} Minimum-loss
  functions, as in \citet[Sections 3.5--3.6]{duartefonseca2026ai}, let the
  researcher ask whether each parameter is pinned down by the chosen moments.
\end{enumerate}

The main weaknesses are also clear. Moment matching throws away information not
contained in the selected moments. It can identify only what the chosen moments
are able to identify. It also introduces a second approximation if \(g(\omega)\)
is learned by a neural moment network. Finally, the choice of weighting matrix
\(W\) and the choice of moments can materially affect the answer.

\subsection{Advantages of likelihood filtering and HMC}

The likelihood route is more ambitious.

\begin{enumerate}[leftmargin=2em]
  \item \textbf{It uses the timing of the data.} A state-space likelihood uses
  the sequence of leverage, investment, valuation, and default observations,
  not just their averages or autocorrelations. Timing is especially valuable
  for default because the model predicts when a firm crosses the default
  boundary.

  \item \textbf{It handles latent states explicitly.} Filtering gives a
  distribution over unobserved capital, productivity, and default risk at each
  date. This can be useful for posterior predictive checks and for studying
  hidden firm quality.

  \item \textbf{It gives a full uncertainty statement.} In the Bayesian version,
  HMC or particle MCMC produces posterior draws of parameters and latent states.
  This gives uncertainty intervals for structural parameters, policies, default
  probabilities, and counterfactuals.

  \item \textbf{It can use irregular and partially missing observations.} A
  state-space likelihood can omit missing components date by date while still
  updating beliefs with the observed components.

  \item \textbf{It tests the model more harshly.} Because the likelihood tries
  to explain the whole observed path, it exposes misspecification that moment
  matching may hide.
\end{enumerate}

The costs are substantial. The researcher must specify a full measurement
model, initial distribution, exit rule, and default-observation process. The
filter can be expensive for large panels. Hard default makes the likelihood
nonsmooth and non-Gaussian. Plain HMC is suitable only for a differentiable
posterior, so the exact default model usually points toward particle MCMC,
while HMC points toward a smoothed approximation whose sensitivity must be
checked. Finally, if the neural solution has errors near the default boundary,
the likelihood can be biased in a very precise-looking way.

\subsection{Side-by-side comparison}

The practical comparison is as follows.

\begin{itemize}[leftmargin=2em]
  \item \textbf{Object being matched.} Moment matching matches selected
  statistics. Filtering matches the probability law of the observed paths.

  \item \textbf{Data requirement.} Moment matching can work with repeated
  cross sections or a panel summarized by moments. Filtering is most valuable
  when firm-level histories and default timing are observed.

  \item \textbf{Modeling burden.} Moment matching needs moment definitions and
  a weighting matrix. Filtering needs a full observation model, latent-state
  transition, initial distribution, and default/exit rule.

  \item \textbf{Computation.} Moment matching is expensive while building the
  simulation and moment-network dataset, but cheap afterward. Filtering is
  expensive at every likelihood evaluation unless the filter is very efficient.

  \item \textbf{Default treatment.} Moment matching can target default
  frequency and related moments without explaining every default date.
  Filtering must assign probabilities to default timing, which is more
  informative but also more fragile.

  \item \textbf{Uncertainty.} Moment matching can use asymptotic standard
  errors, bootstrap, or sensitivity across moment-network folds. Bayesian
  filtering gives posterior uncertainty directly, conditional on the full
  probabilistic specification.

  \item \textbf{Main failure mode.} Moment matching can fit chosen moments while
  missing important dynamics. Filtering can look precise while being wrong
  because the observation model or neural solution is misspecified.
\end{itemize}

\subsection{Which one should we use?}

For a first empirical risky-debt project, moment matching is usually the safer
starting point. It is easier to explain, easier to debug, and less dependent on
strong assumptions about measurement errors and exact firm histories. It is
especially attractive when the goal is to match corporate-finance facts such as
average leverage, investment volatility, issuance frequency, and default rates.

The likelihood-filtering route is attractive when the data contain reliable
firm-level histories and the researcher cares about latent states or default
timing. It is also useful when one wants posterior distributions for
counterfactuals, not just a point estimate. For the risky-debt model, the most
credible likelihood implementation is likely a two-step comparison:

\begin{enumerate}[leftmargin=2em]
  \item Estimate first by moment matching to find a plausible parameter region.
  \item Train or refine the parameter-conditioned neural solution in that
  region.
  \item Run the exact-discrete particle-filter estimator as the benchmark
  likelihood method.
  \item Run smoothed HMC as a faster differentiable approximation.
  \item Check whether the smoothed HMC posterior agrees with the exact
  particle-filter benchmark and whether solver residuals are small throughout
  the posterior region.
\end{enumerate}

This hybrid workflow uses the strengths of both approaches. Moment matching
finds the region of the parameter space that deserves attention. Filtering
then asks the sharper question: can the solved model explain the actual paths
of firms, including the timing of default?

\section{The all-in-one expectation idea}

Maliar et al.'s all-in-one integration idea is easiest to understand through a
small algebra fact. Suppose \(f(\varepsilon)\) is a random number depending on a
shock \(\varepsilon\). Let \(\varepsilon_1\) and \(\varepsilon_2\) be independent
draws from the same distribution.

\begin{proposition}[Two-draw identity]
\label{prop:aio-identity}
If \(\varepsilon_1\) and \(\varepsilon_2\) are independent and identically
distributed, then
\begin{equation}
  \left(\E[f(\varepsilon)]\right)^2
  =
  \E[f(\varepsilon_1)f(\varepsilon_2)].
  \label{eq:aio-identity}
\end{equation}
\end{proposition}

\begin{proof}
Because the two shocks are independent, the expectation of the product equals
the product of expectations:
\[
  \E[f(\varepsilon_1)f(\varepsilon_2)]
  =
  \E[f(\varepsilon_1)]\E[f(\varepsilon_2)].
\]
Because the two shocks have the same distribution,
\[
  \E[f(\varepsilon_1)]=\E[f(\varepsilon_2)]=\E[f(\varepsilon)].
\]
Substituting gives \cref{eq:aio-identity}.
\end{proof}

This identity explains why two independent shock draws can estimate certain
squared expected residuals. Maliar et al. use this idea to reduce the cost of
numerical integration in Bellman and Euler residual objectives
\citep[Definitions 2.9 and 2.10]{maliar2021deep}. In the risky-debt model, this
identity is directly useful only for squared residual components that are linear
in a conditional expectation before squaring. The full risky-debt objective also
contains nonlinear pieces, such as the Fischer-Burmeister expression in
\cref{eq:bellman-fb-residual}, and possibly a smoothed default classifier. For
those pieces, a practical implementation should use mini-batch inner averages or
a separately justified reformulation.

\section{A concrete algorithm}

A practical version of the method is as follows.

\begin{enumerate}[leftmargin=2em]
  \item Choose neural-network architectures for
  \(\widehat V_\theta\), \(\widehat h^k_\theta\), \(\widehat h^b_\theta\), and
  \(\widehat r_\theta\).

  \item Decide which constraints will be hard-coded in the output layer. For
  example, use sigmoid transformations for bounded choices \(k'\) and \(b'\),
  and use softplus transformations for positive quantities such as gross risky
  returns. Keep Bellman/default residuals for the default boundary, because
  default is a regime choice rather than an ordinary output bound.

  \item Draw an initial batch of states \(s_i=(k_i,b_i,z_i)\). These can be
  drawn from a broad rectangle or simulated from a preliminary policy.

  \item For each state, compute the policy
  \((\widehat k'_i,\widehat b'_i)\).

  \item Draw future shocks using
  \[
    \log z'_i=\rho\log z_i+\varepsilon'_i.
  \]

  \item Compute the continuation gate \(C_\theta\), the interior gate
  \(I_\theta\), and the product \(S_\theta=C_\theta I_\theta\).

  \item On states with \(S_\theta\approx 1\), compute the Euler residuals
  \(\widehat{\mathcal E}^k_\theta\) and
  \(\widehat{\mathcal E}^b_\theta\). On states near default or action bounds,
  compute the Bellman/default residual \(\mathcal B_\theta\) and any KKT
  residuals \(\mathcal K_\theta\). If the relevant action bounds are
  hard-coded by output transformations, the corresponding KKT residuals can be
  omitted.

  \item For every positive debt choice \(\widehat b'_i>0\), compute the pricing
  residual \(\mathcal P_\theta\), because the risky rate must make lenders break
  even whether the firm is close to default or safely solvent.

  \item Either minimize the empirical hybrid loss
  \(\mathcal L_N^{\mathrm{hybrid}}(\theta)\) by stochastic gradient descent or
  Adam, or use the hard-constrained hybrid loss
  \(\mathcal L_N^{\mathrm{hard\text{-}hybrid}}(\theta)\) when feasibility is
  enforced by construction. Alternatively, split the same residuals into
  separate objectives as in \cref{subsec:multiobjective-training}, compute one
  gradient per objective, and update the network with a multi-objective
  aggregate such as MGDA. A direct policy-maximization loss
  \(\mathcal L_N(\theta)\) can also be used, but it usually requires sampling
  many alternative actions.

  \item Resimulate states using the updated policy and repeat.

  \item Assess accuracy on fresh states not used in training. Report Bellman
  residuals, Euler residuals by region, KKT residuals near bounds, pricing
  residuals, MOO diagnostics when MOO is used, default classification behavior,
  direct checks against random alternative actions, and simulated financial
  moments.
\end{enumerate}

\section{What this derivation does and does not prove}

The derivation proves a mathematical reformulation: the risky-debt Bellman
problem, the default option, and the zero-profit debt-pricing equation can be
written as residual equations suitable for the deep-learning framework of
\citet{maliar2021deep}. It also shows why the risky-debt case needs more
residuals than the basic Section 3.1 investment model: the risky rate
\(\widetilde r\) must be learned or solved jointly with the value and policy
functions.

The derivation does not prove that a neural network trained with a particular
optimizer will converge to the true economic solution. Neural networks are
approximating families, and stochastic gradient methods are numerical
optimization tools. Multi-objective training can reduce sensitivity to fixed
loss weights, but it is still a numerical optimizer with design choices:
objective grouping, normalization, learning rate, aggregation rule, and stopping
criteria all matter. Euler and KKT residuals are local necessary conditions, so
they should be supplemented by direct checks against alternative actions when
global policy optimality matters. The right scientific claim after training is
therefore not ``the model is solved exactly.'' The right claim is: the trained
functions solve the model approximately on the tested state region, with
approximation quality measured by out-of-sample Bellman, pricing,
Euler-by-region, KKT, direct policy-check, multi-objective, and simulation
diagnostics.

\begin{thebibliography}{9}

\bibitem[D{\'e}sid{\'e}ri(2012)]{desideri2012mgda}
D{\'e}sid{\'e}ri, Jean-Antoine.
\newblock 2012.
\newblock Multiple-gradient descent algorithm ({MGDA}) for multiobjective
optimization.
\newblock \emph{Comptes Rendus Math{\'e}matique} 350(5--6):313--318.
\newblock DOI:
\href{https://doi.org/10.1016/j.crma.2012.03.014}{10.1016/j.crma.2012.03.014}.

\bibitem[Andrieu et~al.(2010)Andrieu, Doucet, and Holenstein]{andrieu2010pmcmc}
Andrieu, Christophe, Arnaud Doucet, and Roman Holenstein.
\newblock 2010.
\newblock Particle Markov chain Monte Carlo methods.
\newblock \emph{Journal of the Royal Statistical Society: Series B}
72(3):269--342.
\newblock DOI:
\href{https://doi.org/10.1111/j.1467-9868.2009.00736.x}{10.1111/j.1467-9868.2009.00736.x}.

\bibitem[Durbin and Koopman(2012)]{durbinkoopman2012ssm}
Durbin, James, and Siem Jan Koopman.
\newblock 2012.
\newblock \emph{Time Series Analysis by State Space Methods}.
\newblock Second edition. Oxford University Press.

\bibitem[Duarte and Fonseca(2026)]{duartefonseca2026ai}
Duarte, Victor, and Julia Fonseca.
\newblock 2026.
\newblock AI for structural estimation.
\newblock NBER Working Paper No. 35283, May 2026.
\newblock Available at:
\href{https://www.nber.org/papers/w35283}{https://www.nber.org/papers/w35283}.

\bibitem[Hall-Hoffarth(2023)]{hallhoffarth2023hard}
Hall-Hoffarth, Emmet.
\newblock 2023.
\newblock Non-linear approximations of DSGE models with neural-networks and
hard-constraints.
\newblock arXiv:2310.13436.
\newblock Available at:
\href{https://arxiv.org/abs/2310.13436}{https://arxiv.org/abs/2310.13436}.

\bibitem[Maliar et~al.(2021)Maliar, Maliar, and Winant]{maliar2021deep}
Maliar, Lilia, Serguei Maliar, and Pablo Winant.
\newblock 2021.
\newblock Deep learning for solving dynamic economic models.
\newblock \emph{Journal of Monetary Economics}.
\newblock DOI: \href{https://doi.org/10.1016/j.jmoneco.2021.07.004}{10.1016/j.jmoneco.2021.07.004}.

\bibitem[Neal(2011)]{neal2011hmc}
Neal, Radford M.
\newblock 2011.
\newblock MCMC using Hamiltonian dynamics.
\newblock In \emph{Handbook of Markov Chain Monte Carlo}, edited by Steve
Brooks, Andrew Gelman, Galin Jones, and Xiao-Li Meng. Chapman and Hall/CRC.
\newblock Available at:
\href{https://www.mcmchandbook.net/HandbookChapter5.pdf}{https://www.mcmchandbook.net/HandbookChapter5.pdf}.

\bibitem[Sener and Koltun(2018)]{sener2018mtlmoo}
Sener, Ozan, and Vladlen Koltun.
\newblock 2018.
\newblock Multi-task learning as multi-objective optimization.
\newblock In \emph{Advances in Neural Information Processing Systems}
(NeurIPS), pages 525--536.
\newblock Available at:
\href{https://proceedings.neurips.cc/paper/2018/hash/432aca3a1e345e339f35a30c8f65edce-Abstract.html}{NeurIPS proceedings}.

\bibitem[Strebulaev and Whited(2012)]{strebulaev2012dynamic}
Strebulaev, Ilya A., and Toni M. Whited.
\newblock 2012.
\newblock Dynamic models and structural estimation in corporate finance.
\newblock Working paper, revised December 2012. This note uses the December
2012 working-paper version for section and equation numbering.
\newblock Available at SSRN:
\href{https://ssrn.com/abstract=2091854}{https://ssrn.com/abstract=2091854}.

\end{thebibliography}

\end{document}
